% Original notes 27 and 28 have plain callouts but starred body labels.
\providecommand{\fnoteplainmark}[3]{%
  \begingroup\renewcommand{\thefootnote}{#1}\footnotemark
  \renewcommand{\thefootnote}{#1#2}\footnotetext{#3}\endgroup
}
\chapter{B)~皮埃尔与Motif}

\section{IV.~Motif——或一次诞生的下葬}

\subsection{一个梦的回忆——或Motif的诞生……}

\notehead{51}{[该注由注46所召]}

\pmark (4月19日)写下这些文字(它们是“我的孤儿们”一注,即n\textsuperscript{o}~46的收尾)还不到一个月,我已发现它们稍稍落在了事件的后面!我刚收到“Hodge Cycles, Motives and Shimura Varieties”(LN~900),作者为Pierre~Deligne、James~S.~Milne、Arthur~Ogus与Kuang-Yen~Shih——那是Deligne惠寄的,还附了一份他的论著目录。\origpage{205}这部汇集六个文本的文集1982年出版,是1970年以来一个有趣的新事实:Motif观念出现在标题里,并在正文中占有了一席之地——虽然还极其谦逊,主要是经由“MotifGalois群”的观念。当然,距离一部Motif理论的全景图还很远;十五二十年来,这幅全景图一直在等待那位勇敢的数学家,肯把它“描绘出来”——足够宽广,以充当一代或几代算术几何学家的灵感之源、引导线索与地平线,他们将有幸确立其有效性(或无论如何,发现Motif之实在的底细……)(注53)。

好像也是从1982年起\fnote{1}{(*)}{(5月25日)我又落后了——这一回是一年:转折发生在1981年6月的吕米尼会议;参见“不义——或归来的意义”一注,n\textsuperscript{o}~75。},时尚的风向对着导出范畴多少开始转了;Zoghman~Mebkhout(在一阵或许有点欣快的高谈里)已经看到它们即将“席卷数学的所有领域”。它们的用处——任何一个消息灵通的人,单凭数学本能,早在六十年代初就已看得很清楚——如今才开始被承认;这(在我看来)主要得归功于Mebkhout孤身的努力:整整七年,他顶着一场专横的时尚,以只信自己本能者的勇气,扛下了为人开路、承受初试代价这桩吃力不讨好的差事……

\pmark 引人注目的是,读这部第一次让(在我退出数学舞台十二年之后)Motif观念低调重返“通行”数学观念之列的出版物,任何不知情的读者都无从疑心:\origpage{206}我这个人,竟与这个长期遭禁的观念的诞生、与围绕它发展起来的一部丰富而精确的“瑜伽”,有过任何关联——在那里,它(以极其残缺的形态)宛如凭空出现,不提任何父系渊源(注51$_{1}$)。

仅仅三周之前,我花了一两页篇幅大谈Motif瑜伽——就连我的一个“孤儿”,而且是我比任何一个都更挂心的那个,也表明我当时显然完全没有说到点子上!毫无疑问,我是在做梦:当我自觉得记起多年孕育的岁月——一种先是微弱而难以捉摸、随后逐月逐年丰富与明确的视野,在那顽强努力的尝试中,试图把握诸种已知的上同调理论(注54)所各自体现的那个“motif”、那共同的本质——它们每一种都以自己的语言,谈论着作为其直接可触的表现之一的那个“motif”的本性。毫无疑问,我还在做梦:当我记起Serre的某个直觉给我留下的强烈印象——Serre被引向把一个投射有限Galois群(一个本质上属于离散性质的——或至少依定义可归结为一个\textbf{有限}群投射系统的——对象)看作这样一个系统的源头:一个庞大的$\ell$-进\textbf{解析}群、乃至$\mathbb{Q}_\ell$上\textbf{代数}群(经由合宜的代数包络)构成的投射系统;这些群甚至往往是约化群——随之,解析群与代数群的(Lie式)直觉与方法的整套武库也被引入。这一构造对每个素数$\ell$都有意义;而我感到(或者我梦见我感到……)这里有一个待探测的奥秘:对不同素数,这些代数群之间的关系如何;它们必定全都出自同一个代数群投射系统,位于所有这些基底域唯一自然的公共子域之上,即特征零的“绝对”之域$\mathbb{Q}$。既然我喜爱做梦,我就继续梦见:我记起自己经由一项肯定只是梦的工作(因为我什么也没有“证明”)进入了这个隐约瞥见的奥秘;我记起自己终于懂得,Motif观念如何提供了理解这一奥秘的钥匙——如何仅仅凭着一个范畴的在场(这里是给定基底概形\dmark{}\origpage{207}上的“光滑”Motif的范畴,比如给定基底域上的Motif),它带有与域$k$上一个代数pro-群的线性表示范畴上所见相似的内蕴结构(代数pro-群观念的魅力,先前也是Serre向我揭示的),人们就真能重建出这样一个pro-群(一旦有了一个合宜的“纤维函子”),并把“抽象”范畴诠释为它的线性表示范畴。

通往一部“MotifGalois理论”的这一进路,是我多年前为描述一个拓扑空间或概形(乃至任意一个topos——不过在这里我感到,自己要刺痛那些“topos并不能使他们开心”的娇贵的耳朵……)的基本群所找到的那条进路向我提示的:即经由所论“空间”的平展覆盖的范畴,以及其上的纤维函子。而“MotifGalois群”的\textbf{语言}本身(我也满可以叫它们Motif“基本群”——两类直觉,自五十年代末以来,对我是同一回事……),连同“纤维函子”的语言(它们与上文所说的那些“显明的化身”——即适用于一给定Motif范畴的各个“上同调理论”——恰相对应),正是为表达这些群的深层本性而造的语言,并且显然提示了它们与Galois群及寻常基本群的直接联系。

我还记得在纤维函子这场游戏里的乐趣与惊叹,还有在Galois群作用下、使不同纤维函子通过“扭转”彼此过渡的torseur——在一场格外具体而迷人的境况里,重新找回Giraud著作中所发展的非交换上同调的整套观念:由纤维函子构成的gerbe(这里的底是平展topos,或者更好,$\mathbb{Q}$的fpqc topos——非平凡而有趣的topos,若曾有过的话!),以及束住这条gerbe的“带”(lien,以群或代数pro-群为值),连同这条带的诸种化身——经由各样代数群或代数pro-群实现,对应于gerbe的各个“截面”,也就是各个上同调函子。比如,特征零概形的各个复点,经由相应的Hodge函子,生成了gerbe的那么多截面,以及从一个截面到另一个截面的torseur;这些torseur连同其上运作的pro-群,带有卓越的代数-几何结构,表达着Hodge上同调的特有结构——不过在这里\dmark{}\origpage{208}我抢在了Motif之梦的另一卷前面……那是这样的年代:今天引领时尚的人们尚未宣布,topos、gerbe之类不逗他们乐,谈这些就是胡扯(当然,就算他们这么说,也丝毫不会妨碍我在它们所在之处认出topos与gerbe……)。而如今十二年又过去了,同一些人却装出发现并教导的样子说:gerbe(即便还不是topos),确实与代数簇的上同调、甚至与阿贝尔积分的周期,大有关涉……

在这里我本可以唤起另一段记忆的梦(或另一个梦的记忆……),它也生于一记“强印象”(我分明深陷于主观性之中!):Serre就Weil猜想背后某种“哲学”所作的若干评论。把这些猜想翻译成$\ell$-进系数($\ell$变动)的上同调语言,令人隐约感到,相应的上同调带有卓越的结构——“权滤过”的结构\fnote{2}{(*)}{(1985年1月24日)对这段走了样的记忆的更正,见注n\textsuperscript{o}~164(I4)及子注n\textsuperscript{o}~164$_{1}$,其中对“权的瑜伽”的渊源有更准确的说明。}。不同$\ell$-进上同调所共有的那个“motif”,想必正是这一本质算术结构的终极载体;由此,这一结构取得了几何的面貌——即几何对象“motif”上的一种卓越结构。把它称作“工作”(尽管这分明仍旧不折不扣是猜谜的片段)肯定又是我在自欺——当需要“猜出”(唯一的向导,是一个正在成形的视野的内在自洽,借助于此处或彼处已知的或猜出的零星要素……)一个Motif的各个“化身”的上同调所特有的结构:权滤过\dmark{}\fnote{3}{(**)}{(1985年2月28日)这里我的头脑有一点点混乱。实际上,说的是与之紧密相连的按“层”(niveaux)的过滤。}在其中如何体现——从Hodge化身开始(在Hodge--Deligne理论尚未问世的时代,而且不无原因……\fnote{4}{(***)}{那是一个年轻的Deligne恐怕还没有在数学语境里听说过“概形”一词、也没听说过“上同调”一词的时候。(这些观念,是他自1965年起在与我的接触中认识的。)})。这使我在(梦中)得以看到,Tate关于代数闭链的猜想(这已是第三记“强印象”,它在Motif之梦里激励着做梦者!)与Hodge猜想(注55)汇入同一幅宏阔的图景,并从中提炼出两三条同一成色的猜想;我对某些人谈起过它们,他们想必已把它们忘了,因为我从此再没听到有人提起,正如再没人提起“标准猜想”一样\dmark{}。无论如何,\origpage{209}它们只是猜想(而且并未发表……)。其中一条并不涉及某个特定的上同调理论,而是给出:一域上非奇异射影簇的Motif上同调上的权滤过,经由该簇自身按给定余维数的闭子集所作的几何过滤的直接诠释(余维数扮演“重量”的角色)\fnote{5}{(*)}{(1985年2月28日)实际上,这里是按“层”(niveaux)的过滤(参见前一条脚注)。}。还有那项“猜出”权在六种运算下的行为的工作(我本该给“工作”打上引号,可又下不了决心!)(那些工作早已失传……)。在那里,我同样从不曾有发明之感,而始终是发现——不如说,是听事物向我言说,当我肯手握笔杆费心去听的时候。它们所说的,有一种不容错辨的精确。

然后有了第三个“Motif之梦”,它像前两个梦的联姻——这一回,是要诠释Motif范畴所带的各种附加结构(其中最首要的,恰恰就是权滤过):用MotifGalois群上的结构,以及以这些群为结构群、用来“扭转”一个纤维函子以(典范地)得到任一其他纤维函子的torseur\dmark{}\fnote{6}{(**)}{正如某个“空间”$X$在两个“点”$x$与$y$处的基本群$\pi_1(x)$、$\pi_1(y)$,经由“扭转”——即经由从$x$到$y$的道路类的群$\pi_1(x,y)$这个torseur——即可彼此约化……}的语言。我记得,在这里更谈不上什么猜谜,而是规规矩矩的数学翻译。它们是一道道未发表的关于代数群线性表示的“练习题”,我兴味盎然地做了几天、几星期,分明感到自己正一步步逼近一个多年来令我心醉神迷的奥秘!必须以表示的语言去把握与表述的观念里最精微的,或许是Motif的“极化”观念——我由Hodge理论取得启发,力图从中澄清在Motif语境中仍有意义的东西。这一反思,应当就发生在我思考“标准猜想”之表述的前后——两种思考\dmark{}\origpage{210}皆受启于Serre(永远是他!)关于Weil猜想的“kählérien”类似物的想法。在这样的境况里,当事物本身向我们低语:它们隐藏的本性为何,我们经由何种手段能最精妙、最忠实地表达它——尽管许多本质事实似乎还远在证明的射程之外——单纯的直觉告诉我们:把事物执拗地向我们低语的一切,径直白纸黑字写下来;我们越肯听写,写下的就越清晰!无须为完备的证明或构造操心——在这一阶段背负这类要求,无异于自绝于一项大面积发现工作中最精微、最本质的阶段:视野的诞生——它从一派表面的虚无中取得形态与实体。单单\textbf{书写}、\textbf{命名}、\textbf{描述}这一行为——哪怕起初只是描述那些难以捉摸的直觉、或那些迟迟不肯成形的“疑点”——就具有一种\textbf{创造性的力量}。当认识的热情投入理智所能把握的事物之中,这正是它手中的头号工具。在这一类事情的发现之路上,这一工作是所有环节中最具创造性的环节:它永远先于证明,并为我们提供证明的手段——或者不如说,没有它,“证明”某事的问题根本无从提出,因为凡是触及本质的东西,还没有一样得到表述与看见。

单凭表述努力的力量,无形的开始成形,变得可以审察;明摆着是错的与可能的得以澄清;而其中那些与已知的或猜出的事物之整体如此完美地协调的,又转而成为正在诞生的视野的坚实可靠的要素。这一视野在表述的工作中日益丰富、日益明确。十样仅仅被猜到的东西——没有一样(比如Hodge猜想)单独令人信服——却相互映照、相互补足,似乎汇入同一种尚带神秘色彩的和谐,便在这和谐中获得视野的力量。即便这十样到头来全被证明是错的,那通向这一临时视野的工作也没有白做;它使我们隐约瞥见、并容我们稍稍深入的那种和谐,不是幻象,而是实在,召唤着我们去认识它。正是经由这一工作,唯有经由这一工作,我们才得以与这一实在、这一隐藏而完美的和谐亲密接触。当我们懂得:事物之为此物自有其理\dmark{}\origpage{211},我们的天职是认识它们而非支配它们,那么,一个错误曝光的日子,就是欢腾的日子(注56)——正如某项证明使我们毫无疑义地确知:我们曾如此想象的某事,当真是实在本身的忠实而真正的表达——那一天同样是欢腾的日子。

在这两种情形里,这样的发现都是一项工作的报偿,没有那工作,发现无从发生。但它或许要待经年累月的努力之后才姗姗而来,或许我们永远听不到那最后的底细——那是留给后来者的——而工作本身就是它自己的报偿,在每一瞬间都因这一瞬间向我们揭示的东西而丰盈。

\notehead{51$_{1}$}{}

(6月5日)然而Zoghman~Mebkhout刚刚提请我注意:前引书第261页,在Deligne一篇“重写并补充一封致Langlands的信”的文章里,提到了“Grothendieck的Motif”。那里写道:“这里要说的不是Grothendieck的Motif——他那样用代数闭链来定义的Motif——而是\textbf{绝对Hodge Motif},同样用绝对Hodge闭链来定义。”这里点名“Grothendieck的Motif”(不加着重),不是作为灵感之源,而是为了与之划清界限,并强调说的是\textbf{另一回事}(人们特意加了着重)。这一撇清格外引人注目,因为Hodge猜想的成立(我假定Deligne,一如其信函式文章的任何读者——从最初的收信人Langlands起——都知道这条猜想)将蕴含这两个观念\textbf{同一}!!

当然,早在1964年——我发展MotifGalois群观念的那一年——我就清楚地知道,可以在同一个模型上发展“Hodge Motif”观念,以及相应的“Galois-Hodge Motif群”观念;后者后来由Tate独立引入(我说不好在他之前还是之后),并得到了“Hodge-Tate群”(联系于一个Hodge结构)之名。这桩粗劣的欺骗(但出自这样一位显赫人物,它似乎并不叫任何人感到不适)在于:以一个可笑的借口——即通往该观念的技术进路(经由绝对Hodge闭链,而非代数闭链)与我所(极临时地)采取的进路不同(倘若Hodge猜想是错的)——把一个新颖而深刻观念的父系渊源干脆整个地藏掖起来,连同我围绕Motif发展起来的一整片丰富的直觉织体。我发展了近十年之久的这部瑜伽\dmark{}\origpage{212},自1968年肇始,就是Deligne著作中灵感的主要来源。它的多产与作为发现工具的力量,早在我1970年离去之前就昭然若揭;而它的同一性,不依赖于为确立它某一部分有限内容之有效而采取的任何技术进路。Deligne有幸提炼出这样两条互不依赖任何猜想的进路。但他却没有那份诚实点明其灵感之源;相反,从1968年起他就极力对所有人遮掩它,以便把它的收益独占为己有,只等1982年(以默许的方式)认领这份信誉。

\subsection{下葬——或新父亲}

\notehead{52}{}

回到Motif之梦,我记得自己还曾把它大声梦见。诚然,梦的工作按其本性是孤独的劳作——但这项顽强地持续了多年的工作(在一项吞噬了我大部分时间的浩繁奠基写作的边缘上进行)的种种波折,有一位日复一日的见证者,比Serre亲近得多——Serre只满足于远远地看这些事……\fnote{7}{(*)}{(5月25日)不过,我对Motif的反思,其实早于Deligne的出现。我关于MotifGalois理论的手写笔记,标注的日期是1964年。}。关于这位日复一日的心腹之交,我在回顾中写道:他在六十年代中期“多少扮演过学生的角色”,而我“给他讲过我在代数几何里知道的那一点点”。我本可以补上一句:我连按这个词的通常意义我所不曾“知道”的东西也讲给他听——那些数学的“梦”(以Motif为主题,也以别的为主题),它们总能在他那里找到一只专注的耳朵和一颗警醒的心灵,和我一样渴求理解。

诚然,当我写道Pierre~Deligne曾“多少扮演过学生的角色”时,这仍旧全然是一种主观的印象(注57),据我所知,没有任何书面的、至少是印刷的痕迹可资佐证——那种痕迹本可以让随便哪个人疑心,Deligne从我嘴里学到过什么——而在这里,我却乐于记起:我同他谈数学,从来没有一次不从中学到点什么。(甚至在我同他不再谈数学之后,我仍继续经由他学习那些更困难、或许也更重要的事情,包括在我写下这些文字的今天……)。

\pmark 新近,一位第三者(他不知怎么猜到了这事或许会让我感兴趣!)告知我,Deligne和别人有一部文本,其中谈到Motif,或者至少谈到“塔纳卡范畴”;我向Deligne提了一句,他由衷地惊讶,我居然会对这类东西感兴趣。\origpage{213}翻阅他惠寄的那一册,我确信他的惊讶完全有理。显然,我这个人,同书中所论的主题全然无涉。至多在导言里一句话带过:某些“标准猜想”(当年是我提出的,天知道为什么)对一域上Motif范畴的结构会有一个推论……好奇的读者想多知道些就为难了:他在这整部书里,找不到关于这些猜想的任何说明或出处——它们此后就再没被提起;也找不到对唯一一篇我解释如何以标准猜想在域上构造Motif范畴的已刊文本的提及;也找不到1970年以前唯一另一篇谈到Motif的已刊文本——出自Demazure之手(如果我没记错,是在一场Bourbaki研讨班上),它遵循我那条特设的构造原则,只是视角稍有不同……\fnote{8}{(*)}{核实之后,我发现,除了关于标准猜想的几页(Algebraic Geometry, Bombay, 1968, Oxford Univ.~Press (1969), 第193--199页),我没有任何已刊的数学文本谈到Motif。在Demazure的报告(Bourbaki研讨班n\textsuperscript{o}~365,1969/70,继Manin的俄文报告之后)里,提到我1967年在IHES所作的几次报告,它们(我猜想)应当构成Motif 视野的第一幅总草图。关于标准猜想及其与Weil猜想的关系,有一篇比孟买会议通告更详细的报告,出自Kleiman之手(Algebraic Cycles and the Weil conjectures, 载 Dix exposés sur la cohomologie des schémas, Masson-North Holland, 1968, 第359--386页)。在1970年以前,除我之外,我没有见到任何关于标准猜想的反思——尤其朝着对它们的证明。存心无视这些关键猜想(我在孟买的草图里说过,我认为它们连同优秀概形的奇点消解,是代数几何里最重要的未决问题)的做法,在我看来,很大程度上造成了代数簇上同调理论给我的那种停滞印象——就传到我耳中的种种回响而言。}

\pmark 连Neantro~Saavedra——他有幸位列我“1970年以前的学生”——都得到了应有的引用\origpage{214}。他跟我做过一篇学位论文,论的是我记得我称之为“刚性张量范畴”、而他称之为“塔纳卡范畴”的东西。人们至今仍在纳闷:Saavedra是靠了怎样奇迹般的运气,恰好预见了十年后才绽开的DeligneMotif理论的需要!实际上,在他的论文里,他做的恰恰就是\textbf{那}项在技术上构成MotifGalois理论之钥匙的工作,一如 J.L.~Verdier的论文原则上就是\textbf{那}项在技术上构成上同调六运算形式主义之钥匙的工作。有一个(此外还有别的)对Saavedra有利的不同:他费心发表了自己的工作;他确实不曾有Hartshorne、Deligne与Illusie联袂的笔杆,来替他免除这道手续。然而十年之后,Saavedra的论文被从头到尾、几乎一字不易地复制进那部卓越的文集里,这一次执笔的是Deligne与Milne。倘若只是为了更正Saavedra工作中的两个具体之点(注58),这或许并非必要。但万事必有其所以然,我想我辨出了Deligne亲自不辞这份辛劳的缘由\fnote{9}{(*)}{对此可参见“白板”一注中的反思,n\textsuperscript{o}~67。}\fnote{10}{(**)}{(6月8日)而当事情涉及带有我的影响之痕迹的工作时,就更是如此——对此可参见“注记——或新伦理”那一情节,第33节。}:那分明同他自己那条在出版事宜上把严格推到极点的标准背道而驰;而众所周知,一涉及别人,他把这条标准执行得一丝不苟……

至于诸观念与Motif瑜伽本身的父系渊源,对不知情的读者(而知情的读者已开始变得稀少,终将一个个寿终正寝……)来说,这一父系渊源无可置疑——无须在这里去惊动遥远的Hilbert与Riemann,更不必去惊动上帝。倘若那位显赫的作者——他关于阿贝尔簇上绝对Hodge闭链的漂亮结果,俨然是Motif理论的起点,干脆说,是Motif理论的诞生——对其父系渊源只字不提,那是一种令他生辉的谦虚\dmark{}\origpage{215},与行业的惯例与伦理完美合拍:按后者的要求,应当把荣耀归于它分明所当得之处的差事(如有必要的话),留给别人:留给合法的父亲……

\notehead{53}{}

被这个孤儿的坎坷际遇所触动,并疑心不会再有别人去做这项工作——这项工作,至今显然只有我一个人感到其需要与规模——我料想,那位“勇敢的数学家”不是别人,正是我自己,待我走完《追寻栈》之路之后(我预计它还要占去我大约一年)。

\notehead{54}{}

从那以后,代数簇又出现了两种新的上同调理论(除Hodge--Deligne理论之外——后者是Hodge上同调在“Motif”精神里的自然延伸),即Deligne的“分层pro-模”理论,以及尤其是晶体理论——Sato--Mebkhout式的“D-模”版本,配上前文说过的“好上帝定理”(又名Mebkhout定理)所带来的新的照亮。这条通往可构造离散系数的进路,多半注定要取代Deligne先前的版本,因为它无疑更适合表达与De~Rham上同调的关系。况且,这些新理论并不在给定概形上光滑Motif的范畴上提供新的纤维函子,而是(尚需一项比迄今所做更深彻的奠基工作)提供一种方式,以精确把握复数域上有限型概形上一个Motif(不必光滑)的“Hodge”化身,或特征零域上有限型概形上的“De~Rham”化身。此外很有可能:$\mathbb{C}$上有限型概形的Hodge--Deligne系数理论(显然始终未见成文),终将显现为被包含在Sato--Mebkhout式晶体系数理论(同样未见成文;并另附一项关键的滤过数据)之中,或者更确切地说,显现为后者与$\mathbb{Q}$-向量空间的可构造离散系数理论的某种交……至于澄清Mebkhout式晶体理论与Berthelot及其他人在正特征下发展的理论之间的关系,那是Mebkhout早在1978年之前就已感到的任务,处在一种普遍冷漠的气候里;而在我看来,它是近期摆在我们面前、关乎对代数簇“那一个”上同调(唯一而不可分,即Motif的上同调!)之理解的最迷人的任务之一。

\notehead{55}{}

\dmark{}我尽管做梦,但我关于Motif与Hodge结构之关系的梦,却让我在不经意间把手指按在了Hodge“广义”猜想——按Hodge最初表述的样子——的一处不自洽上,并把它代之以一个更正的版本;这一回(我敢打赌),它应当既不会比代数闭链上那个“寻常的”Hodge猜想更错,也不会比它少错。\origpage{216}

\subsection{一场大屠杀的前奏}

\notehead{56}{}

我尤其想到——恰恰在代数簇上同调的语境里——Griffiths对一种诱人的想法之错误性的发现:长期以来人们以为,一个同调等价于零的闭链,会有一个代数等价于零的倍数。对一个全新现象的这一发现,当时给我的触动之深,竟使我花了整整一星期的功夫去力求吃透Griffiths的例子,把他的构造(那是超越的,在$\mathbb{C}$上)搬移成一种“尽可能一般”的构造,并且在任意特征的域上成立。这一推广并非一目了然,靠的是(如果我没记错)Leray谱序列与Lefschetz定理。

(6月16日)这一反思,曾是我有机会在平展语境里发展“Lefschetz束”的上同调理论的机缘。我这方面的笔记,已由 N.~Katz在SGA~7~II研讨班(由 P.~Deligne与 N.~Katz主持)的报告XVII、XVIII、XX里详细阐发(他费心指涉了这些笔记,而且紧随其步伐)。而在 P.~Deligne所写的该卷导言里——其中说,该卷的关键结果是报告XV(平展上同调中的Picard--Lefschetz公式)与XVIII(Lefschetz束理论)——作者却小心地不提:在这套Lefschetz束的“关键理论”里,我也有我的一份。读这篇导言,得到的印象是:对该卷所发挥的主题,我毫无贡献。

SGA~7那场漫长的研讨班,1967--69年间接续了1960至1967年间在我推动下发展的SGA~1至SGA~6诸研讨班,曾由Deligne与我共同主持,是我以消没闭链群的一个系统理论发令起跑的。由于各报告由不同的志愿者执笔、拖沓经年,研讨班的两个分卷(SGA~7~I与SGA~7~II)迟至1973年才出版,经手的是\dmark{}\origpage{217}Deligne。尽管在研讨班当时,大家已商定它将以一场共同主持的研讨班的面目出现;我离去之后,Deligne却向我表示了他那个(在我看来很奇怪的)愿望:把研讨班\textbf{一分为二},一部分I以我主持的名义出现,另一部分以他和Katz的名义。我如今在其中辨出一场“战役”,它预演了“SGA~4½之役”——那场战役的目标(除别的以外)是要让SGA~1至SGA~7整个奠基丛书显出这样一副面孔:仿佛它是一套随手采撷的文本汇编,我这个人只在其中扮演一个偶发的、乃至多余的角色——尽管就其精神与构思而言,这套丛书一如EGA《代数几何原理》丛书,是与我个人密不可分的。这一倾向在SGA~4½分卷里表现得非常清楚,甚至粗暴;尤其表现在对SGA~5研讨班的屠杀里——该分卷与那场屠杀密不可分。对此可参见(除别的以外)“白板”与“屠杀”两注,n\textsuperscript{o}~67与87,尤其是“遗骸……”(n\textsuperscript{o}~88)。

(6月17日)SGA~7研讨班的总体构思(我在其中绝不曾区分什么“I”“II”两部分,至今也仍不区分)出自我这里;另一方面,Deligne作出了重要贡献(在我1969年写的Deligne工作评述报告里有记载,见该报告n\textsuperscript{o}~13、14条),其中就研讨班的需要而言最关键的,是Picard--Lefschetz公式——经由从一个已知的超越情形做特化论证而证明。把研讨班切成两部分,无论在数学上还是就双方的贡献而言,都毫无道理——SGA~7的两个“块块”里,Deligne和我都有实质性的贡献。

当然,倘若Deligne继续我开创的SGA奠基丛书——它远未跑完自己的赛程!——我会喜出望外。这场“SGA~7之役”根本不是什么延续;我感到它是一记粗暴的“锯子”(或电锯……),\textbf{终结}了SGA丛书——用的正是一个摆明要与我个人划清界限的分卷,尽管它维系于我的著作,并且像其余各卷一样带着我的印记。当我的个人痕迹在其中被尽量藏掖的同时,对我著作的口吻,尚未达到“SGA~4½之役”那种几乎不加掩饰的轻蔑——后者是对SGA~4与SGA~5研讨班之统一体更深重的一记锯痕,也是对其中未刊的SGA~5部分大加劫掠的手段与借口;劫下的碎块,在Deligne与Verdier之间均分……

\notehead{57}{}

\dmark{}请允许我赶紧补充:同样的评语,也适用于另一位大手笔的数学家——十年之后,我斗胆说过(在注n\textsuperscript{o}~19里)他也“多少扮演过学生的角色”。 \origpage{218}

\notehead{58}{}

这让我想起:Springer的《Lecture Notes》(它曾发表过六七篇在我指导下完成的“1970年以前”的博士论文)从未肯发表Yves~Ladegaillerie“1970年以后”的那篇论文(理由:他们不发表博士论文!)。人们倒可以说,它把Saavedra的论文发表了第二遍……我还曾向Deligne提起Ladegaillerie那条四处碰壁的漂亮同位结果(心里暗暗指望他会出手相助发表)——却没能引起他的兴趣(理由:他对曲面拓扑不在行……)。

幕落……

\subsection{新伦理(2)——或抢夺的集市}

\notehead{59}{}

(4月20日)写下这些文字——它们记录了一桩矛盾及其代价——已有几个星期;我惊讶地发现,当事人早在两年前就已找到一条最简单的路子来“化解”那桩矛盾——难就难在想得到!人们可以称之为“预先下葬法”(读者可在昨天写成的双注(注50、51)里了解它,那是发现之后情绪未凉之时写的)。我感到遗憾的是:这位被预先下葬的死者,冷不防又重新出现在那有名的“数学舞台”上(它有时候,分明更像一个抢夺的集市……),恐怕会给这项光辉方法的干净利落的应用,带来技术上的麻烦!在前面的注里(“共识性的操守——与信息控制”,n\textsuperscript{o}~6),我(还有些含混地)感到:科学行业里那条最获普遍认同的操守规则,在掌管科学信息控制的人们不尊重“每个科学工作者有权让自己的观念与结果为人所知”这一权利时,“就形同虚设”。正是在反思的这个节点上,我还费心相当详细地描述了一个个案:对这一权利的藐视,在我眼里昭然若揭;而且我分明感到,这藐视已逼近对那条首要规则——获得普遍共识的那一条——的藐视。(见“注记——或新伦理”,第30节。)

\pmark \origpage{219}这已不是我头一回感到这种特殊的痛楚:当我看到这条首要规则的\textbf{精神}遭人藐视,而藐视者却处于“暂停”的豁免状态——无论靠的是其地位(高到不容置疑!)、其才具,还是形式上的漫不经心。我试着在相关于所引第30节的那则注里(“年轻人的势利——或纯洁的捍卫者”)勾勒这种痛楚。当人们可以放胆藐视我在那里谈到的那些“显而易见”的事,并且(如今我可以补上一句)以同样的精神,也藐视那些既未证明、也未以已刊并为众人所知的“猜想”的形式注册备案的事情(或许是深刻的事情),那么人们同样可以(鉴于其“分量之轻”\fnote{11}{(*)}{“好上帝定理”(又名Mebkhout定理)的遭遇,正是如此。(6月8日)而且,一如对Motif瑜伽那样,还费心巧妙地制造出对此物拥有父权的表象,却从不明言!就个案而言,可参见“变戏法的人”一注,n\textsuperscript{o}~75$''$;就这条光辉方法的一般路数或风格而言,可参见“暂停!”一注,n\textsuperscript{o}~77,以及紧随其后的“据为己有与轻蔑”一注,n\textsuperscript{o}~59$'$。})把它们视作公共财产(不言而喻,是平凡的),因而,到时候,以最大的漫不经心与最坦然的良心,把它们当作“自己的”——当然,前提是没有人想到要去据为己有一篇十页或一百页(或哪怕十行)的硬核证明——它确立一个“自己本来不会证明”的结果(注59$'$)。我不曾想到,自己对“形同虚设”这一点,竟感触得这样准、说得这样中,因为我有幸目睹:上面那个案子里那条模糊的“界限”,被人轻松越过——而且可以肯定,越过时同样怀着最坦然的良心,\textbf{鉴于其分量之轻:那是一个梦},而且甚至未经证明(尤其,未经\textbf{发表}……)\fnote{12}{(**)}{人们何必有所顾忌,既然事情本身分明表明:如今普遍的共识认为这完全正常——至少出自一位格调如此之高的人物!所谓“坦然的良心”,不多不少,正是这样一种感觉:觉得自己与自己所处行业里通行的共识相合。}。

幸好我还有点招架之力——必要的时候,我多少能表达出我的感受与想说的话;我(不论对错)挣得了一份信誉,因而,当我有话要说的时候,有一线希望被倾听,或者,若我觉得有必要,把它发表出来。相反,我此刻更痛切地体会到那种“不义与无力之感”——当一个人\dmark{}\origpage{220}受了损害却无处申诉,在“一切尽在掌握”的人们的专断面前,自觉手脚都被捆住——而后者还随心所欲地支配一切。

诚然,在我做数学家的生涯里,我也曾以同样坦然的良心有过可鄙的行为;在反思中我已有机会谈到那些个案——它们是这场反思从遗忘与从未检视过的暧昧的迷雾里重新唤起的。探查它们,我终于懂得:今天(而且由来已久)学生轻松愉快地超过老师,我无须惊讶;凡是与我有感情或情谊相连的人,我也无须否认。但对我、对所有的人来说,把猫就叫作猫,是健康的——无论这只猫出自我家,还是出自别人家。

\subsection{据为己有与轻蔑}

\notehead{59$'$}{}

(6月8日)对这件事,就我的朋友Pierre~Deligne而言,我已完全不这么确信了,因为我有机会确认:他终于滑进了对$\ell$-进上同调工具的“默认父权”游戏——也就是我所说的对平展上同调的“掌控”。从“SGA~4½之役”(其中我的名字还被提起,只是带着一种做作的、漫不经心的轻蔑,轻蔑于我著作中的这一核心部分——他自己的那一部分正源出于此),到“葬礼悼词”(其中凡“上同调”一词本身,凡与我的名字相关处,一概禁用),有一段令人瞩目的演变。(初期阶段见“白板”与“格格不入者”两注,终局阶段见“葬礼悼词(1)、(2)”两注。)

作为这场步步升级中的几个中间阶段:1981年,有那篇“值得纪念的文章”,论的是所谓“反常”层(对此可参见“不义——或归来的意义”与“暂停!”两注,n\textsuperscript{o}~75与77);次年,有LN~900里Motif的开棺出土(而“葬礼悼词”则再下一年,即1983年登场)。在我得以观察到的所有这些情形以及一些规模较小的情形里,那种使Deligne能够以完全坦然的良心据有他人观念之信誉的内在态度与“方法”,就是\textbf{轻蔑}(它部分地保持缄默,却又被巧妙地暗示出来)——对着人们即将据为己有的那个“轻”东西:它“轻”到甚至不值一提,尽管人们转过头去就要用它做出真正有力的东西——Weil猜想,所谓“反常”层的理论……一旦战役完成,据为己有成为既成事实、并为众人所接受,总有的是时间来修正口径,带着已据为己有的东西,谦虚地招摇过市。同一项贡献,只要似乎还沾着一个有待下葬者的名字,就遭到漫不经心的轻蔑;\dmark{}\origpage{221}一旦被他本人据为己有,就被捧上云端($\ell$-进上同调,Motif,Mebkhout的瑜伽尚待其时),或者被某位好哥们据为己有(导出范畴的瑜伽,对偶的瑜伽——由Verdier据有,Deligne积极怂恿)。

\section{V.~我的朋友皮埃尔}

\subsection{孩子}

\notehead{60}{}

\dmark{}(4月21日)再拾起这个关于回忆的梦——它不只是对一个视野之诞生的回忆……\origpage{223}我清楚地记得(尽管我忘掉了那么多事情!)那种每一次都焕然一新的乐趣:同那个很快就不止是心腹之交的人交谈——对于一切令我困惑的东西,对于日复一日在我与数学的恋情中豁亮起来、令我心醉的东西,他比任何“学生”都更是知交。他那份始终警醒的兴趣,他了解任何事情的从容(“仿佛他一直就知道……”),对我而言,是一泓永不枯竭的欢喜之泉。他的倾听是完美的,由那种也在我身上跳动着的、求理解的渴望所推动——一种高度醒着的倾听,是心心相印的标志。他的评点,总是迎着我自己的直觉或保留而来;即便并非如此,它们也往往在我力图穿透重重迷雾去把握的那个实在上,投下一道意想不到的光。如我在别处说过的,我提出的问题,他往往应声就有答案,或者在其后的几天、几星期里把答案发展出来。这就是说,倾听是双向的:当他反过来向我解释他找到的答案时——也就是解释事情的所以然,而那所以然,总是带着那种完美的自然、那种从容显现;同样的从容,曾常令我心醉地见于我的某些长辈,如Schwartz与Serre(同样也见于Cartier)。正是这同一种单纯,这同一种“自明性”,是我在理解数学事物的过程中始终追求的。无须明言:凭着这条进路与这份要求,他和我显然“同属一家”。

从我们相识之初,我就分明感到:他那(人们所谓的)“天资”,有一种极其少见的品质,远在我所拥有的那份寒微的手段之上;而就求理解的热情、就对待数学事物之理解的那份严格而言,我们却在同一个音高上。我还隐约感到——当时我尚无法把它表述出来——我在他身上觉察到的这股“力量”(我在自己身上也感到,只是较弱):那种“看见”人人视而不见的自明之事的力量,乃是童年的力量,是孩童眼里的天真。他身上有某种属于孩子的东西,比我认识的任何别的数学家都更外露;这绝非偶然。他给我讲过:有一天,我记得还在中学,他曾兴致勃勃地用定义去验算乘法表(顺水推舟,连加法表也一并验了)\dmark{}\origpage{224},对1到9这些数。他当然不指望有什么意外——如果说真有意外(愉快的,一如既往……),那就是:这证明居然可以在短短几页里、漂漂亮亮又完完整整地做成,前后不过半个钟头的事。当他笑着给我讲这件事时,我分明感到,那半个钟头用得其所——这件事,我今天比当时懂得更透。这段小故事当时打动了我,甚至震动了我(我想我未曾形于颜色)——我从中感到一种\textbf{内在的自主}、一种相对于既有知识的自由的征兆;这种自由,在我童年与数学的关系中,自最初接触起就已在场(注69)\fnote{13}{(*)}{而且我觉得,这份自由在我做数学家的一生中从未全然隐没,如今又像童年时那样在场了。两三年前,我向我的朋友重提乘法表那桩童年小事。我感到他为这一回忆的唤起而局促——它分明已不再符合他心中的自我形象。这份局促我并不真正惊讶,却仍然难过:又一件我明知道、却仍难于承认的事,再次得到了确认……}。这种彼此互为特权对话者的关系——我想,那会儿我们几乎天天见面\fnote{14}{(**)}{至少在我住在比雷时是如此,他当时住在IHES的一间单间。自1967年(我搬到马西)之后,我想我们仍一周会面一两次,至少直到我不再投身数学为止。}——持续了五年,从1965年(如果我没记错)到1969年(含)。我还记得那一年写一份详尽的评述报告时的乐趣:当时我提议把他增选为我自建校起(1958年)就一直工作的那个机构——我的大部分数学著作就是在那里完成的——的教授。那份报告我没有存本了(注64):我在其中评阅了我想有十几篇我朋友的著作,当时几乎都未发表(有许多后来也没发表),而在我看,其中大部分——如果不是全部——都当得起一篇优秀国家博士论文的主体分量。呈递这份雄辩的报告,我感到比呈递一份关于我自己工作的报告更骄傲、更欢喜(那种事,我一生只做过两次,而且每次都是勉强自己……)。其中许多著作,是对我提出的问题的回答(其中唯一已刊的,是前面提到的那篇关于概形真光滑态射之Leray谱序列退化的工作(注63))。而最重要的两篇\dmark{}\origpage{225}却是对Deligne自己提出的问题的回答;显然,其分量与一篇“优秀的国家博士论文”不可同日而语。那就是他关于Ramanujan猜想的工作(原书拼作Ramanuyam)(发表于Bourbaki研讨班),以及关于混合Hodge结构的工作——亦称“Hodge--Deligne理论”。

有一件事很奇怪,而且是我写这份流光溢彩的报告时远远没有料到的:不到一年之后,我就要离开那个机构——我正准备在那里增选我年轻而令人惊叹的朋友,并且打算在那里终老。还有(如今我把这两桩双生的情节放在一起想)另一件怪事,也绝不像是单纯的“偶然”:一两个月前,正是这位(如今已不年轻!)朋友,向我宣布了他自己要从同一机构离去——而恰好一年之前,我才重新拾起一种规律的数学活动,以一种不期然的“重返”走上数学舞台(虽还不是“上流社会”……)

不止一次,我在《收获与播种》里有机会谈到我的离去——谈那次“有益的挣脱”——以及紧随其后的“觉醒”;正是后者使那一情节成为我生命中的一个关键转折。在其后那些紧张的年月里,数学家的世界,连同我在其中爱过的人们,以及数学本身最令我心醉神迷的那个东西,都变得极其遥远——仿佛淹没在另一个“自我”的回忆的迷雾里,那个自我仿佛死去千年……

但无论在这桩情节之前,还是在这第一个大转折之后的年月里,我都知道:那个曾是(一点点\fnote{15}{(*)}{至于我为何有这份顾忌,不肯把(太!)耀眼的Deligne当作我的一个学生,见“格格不入者”一注(n\textsuperscript{o}~67)。})我的学生、(在很大程度上的)心腹之交与朋友的人,只须顺着体内一个游戏着的、渴求认识的孩子那股自发的冲劲,就能发现并唤起崭新而出人意表的世界,就能探测它们、认识它们内在的本性——并由此把它们向他自己的同类、也向他自己揭示。所以,在我离去之后(绝无回归之念!),倘若我看见“一位勇敢的数学家”、一位富有灵感的人,开始大笔勾勒(权作开端……)我隐约瞥见的那幅宏阔图景——而我尚且只画下了一组局部的、临时的草图——那个人只能是他:他把做成这件事的一切都握在手里!勾勒这幅大面积的第一图景——一位“总设计师”把关于\dmark{}\origpage{226}代数簇上同调的已知与猜出的本质事物,纳入一个共同的视野——对于那个心中视野已然呼之欲出、只待破除“尚未写下”之迷雾的人,是几个月的工作,甚至用不上几年。(不妨在此后经年、甚至经代——倘若需要几代的话——将它重做、深化,直到Motif之实在的底细被充分理解与确立。)而且我毫不怀疑,这项不久前还“在我手里灼烧”的工作,随时都会有人去做,至少在其后两三年内,趁着它还热乎。我离去之后,的确只剩一个人,凭着他对认识的冲劲本身,注定要做这件灼热而迷人的工作。不妨,在“总设计师”写就并经检验、大厦的营造或深或浅地推进之后,把继续这项工程的差事——无论它多么迷人——留给别人,自己投身别的冒险,在那个数学事物的世界里:那里的每一个转弯,都显露着一个崭新而无边世界的许诺,只要我们有睁开的新眼睛去看……

当我的生活还在那座温暖的科学暖房里流转——它把尘世的喧嚣隔绝在外——而当Deligne发展他对Hodge理论的推广时(那应当是1968或1969年),在我们之间,有件事不言而喻:那项工作,是把“Motif图景”的某个\textbf{部分}加以实现、检验与明确的一个头等步骤;而那幅图景,作为一个整体,从未被白纸黑字写下\fnote{16}{(*)}{至于那套Hodge--Deligne理论后来(据我所知)从未超出这第一次尝试的阶段,从未扩展成一套“Hodge--Deligne系数”(及其上的“六种运算”)理论、架在复数域上有限型概形之上——这与另一件怪事密不可分:那幅宏阔的“Motif图景”从未被勾勒出来,而它的存在本身,直到今天还被小心地讳莫如深……}。在我告别暖房之后的年月里,当数学对我已相当遥远时,我得知Weil猜想最终获得证明,毫不意外。(若说有意外,那是“标准猜想”没有在同一步里得到证明——尽管人们提出它们,恰恰就是为了逼近Weil猜想,同时也是作为在一域上至少建立半单Motif理论的一个手段\fnote{17}{(**)}{只是在这最后几年里,我才模糊地(但近来更精确地!)意识到:“标准猜想”,连同它们为之提供了第一条“构造性”进路的Motif观念本身,都被\textbf{埋葬}了——缘由如今在我眼前显得格外清楚。(亦可与前面一条脚注对读。)}。)我很清楚:无论是这第一击\dmark{}\origpage{227}——那初试身手的、Hodge式的一般系数理论——还是对若干关键猜想(在若干其他多少为人所知的猜想之中)的这项证明,都还没有让他使出全副本领——差得还远。而我并不着急,因为我的注意力的绝大部分,已倾注在别处。(→注61)

\subsection{下葬}

\notehead{61}{}

我有幸亲眼见到一股童年的冲劲初次绽放,带着大面积舒展的许诺。在其后的十五年里,我渐渐确认:这一许诺始终被不断推迟。他身上有那样一种精微的东西,我曾懂得去感觉、去辨认(尽管那会儿我对那么多事情都无感!)——一种与脑力全然异质的东西(脑力既碾压人也穿透人……),一件对一切真正的创造性的工作而言头等本质的东西。这种东西,我在别人身上有时也感到过;但在我认识的任何数学家身上,它都未曾以可与之相比的力量显现。而我(想当然地)指望,这东西会继续在他身上绽放、转化,并毫不费力地经由一部独一无二的作品表达出来——在那部作品里,我不过是一个寒微的先行者。然而又一桩怪事(而在这么多“怪事”之间,必定有一条简单而深刻的纽带):我眼看这“精微的东西”、这既非膂力亦非脑力的“力量”,随年月流逝而渐渐消退,仿佛被\textbf{埋葬}在一层层越来越厚的沉积之下——那是一层层\textbf{别的东西}的沉积,我对它们知之甚稔——世上最寻常的东西!它未必与脑力、与老到的经验或在某一学科里磨炼出的嗅觉不能相安;这些东西,凭其作品的积累——也许光灿夺目,肯定自有其力量与美——可以博得一些人的叹服、另一些人的畏惧,或兼而有之。但这并不是我当初说“舒展”或“绽放”时所想的。\textbf{那个}绽放,我所想的那一个,是天真的果实——渴望认识,并且随时准备为这个取之不竭的世界里或大或小的事物的美而欢喜,或为它的某个部分(比如数学事物的那个宏阔世界……)的美而欢喜。唯有它,才有深刻革新的力量:无论是自我的革新,还是对世上事物之认识的革新。在我看来,把这一点实现得圆满的,正是谦逊的Riemann\fnote{18}{(*)}{Riemann(1826--1866)的著作,装满区区一册薄薄的文集,约十篇论文(诚然,他四十来岁就去世了),其中大部分包含一些简单而本质的观念,深刻革新了他那个时代的数学。}。这份真正的绽放,与轻蔑无缘\dmark{}\origpage{228}:对他人(对那些自感远在自己之下的人……)的轻蔑,对那些太“小”、太自明以致不屑一顾之事的轻蔑,对那些被认定低于自己正当期望之事的轻蔑;或者还有对某个\textbf{梦}的轻蔑——那梦执拗地向我们诉说着我们口口声声说爱的事物……它与轻蔑无缘,正如它与滋养轻蔑的自负无缘。

诚然,凭着那令人叹服的“天资”,但更凭着那件不打动任何人、却\textbf{创造}着的精微之物,“学生”注定要远远超过“老师”。我毫不怀疑,在我离开那个地方——我在那里曾见证一次如此美丽的起飞——之后的年月里,Deligne将使一部宏阔而深刻的作品得以全副展开,而我曾是它的先行者之一。这样一部作品的回响,岁月流转,不会不传到我这里;而我自己,在远离数学的其他追寻中跋涉,对他将要发现的那些崭新世界的全部分量与全部美,只能有不太完整的领会。

但学生若在内心深处\textbf{否认}老师,努力在人前人后抹去老师的一切痕迹(无论这份馈赠是为了更好还是为了更坏……),学生就无法真正超过老师——正如儿子若否认父亲,就无法真正超过父亲。这一课,我主要是经由我与我的孩子们的关系学到的,但后来也经由我与昔日某些学生的关系;尤其是与那一个人——在所有人之中,我始终有所顾忌、不肯以“学生”之名相称的那一个;因为从相识那一刻起,我就分明感到,我有东西要向他学习——一如他也有东西要向我学\fnote{19}{(*)}{(6月14日)关于我身上这份执拗的蓄意——贬低我所能给的东西、否认一段师生关系的实在——见“格格不入者”一注,n\textsuperscript{o}~67$'$。明摆着:我的朋友从我的接触中所学到的东西(“仿佛他一直就知道”,当然!),与我经由他所学到的,没有共同的尺度。倘若我继续高强度的数学投资直到今天、我们之间保持着规律的数学接触,情形想必会是两样。}

但只是在相识将近十年之后——在1975年之后,尤其自从我开始沉思我所经历的、我所目睹之事的意味——我才开始感到那个继续为我所珍的人身上的这份\textbf{枷锁}。我也隐约感到:这份对我这个人、对我在其人生关键年月里所扮演角色的秘密否认,更深地,也是对他\textbf{自己}的否认。(想必每当我们否认、并想抹去\dmark{}\origpage{229}某件分明发生过、而其果实有待我们去摘取的事情之时,都是如此。)

然而,由于对“数学界时下在做什么”、对他自己在其中所做的事情\fnote{20}{(*)}{自1970年以来,我收到过Deligne的四份抽印本,都当即匆匆翻过(一如我至今仍不时收到的大多数抽印本)。凭这点东西,难以对一部数学著作——哪怕只就其总轮廓或主要主题——形成概念。}多少不“在局中”,在几个星期前的反思之前,我从不曾掂量:这份枷锁,对他倾其所有投注的事情——他的数学著作——也压得有多重。诚然,八九年来,我不止一次看到,一份蓄意的轻蔑口吻(对我)或鄙夷(对一些他有能力使之气馁的别人)仿佛把单纯的常识、把数学家的健全本能一笔勾销(注66)。而且他并非我旧学生中——带或不带引号——唯一一个让我目睹这类态度对着我所珍视的人(或对着别人)的。但在别人身上,它从不曾这样痛切地触动我。在刚过去的两个月的反思里,我不止一次影射过这份经验——“它是我做数学家的一生中所曾经历的最苦的经验”;我也说过,这场《收获与播种》的反思到头来就它教给了我什么。那份创痛如此剧烈,它就一个始终为我所珍的人教给我分量如此之重的东西(而我继续回避它同时就我自己、就我的过去所教给我的……),以至于这种态度会对他、乃至对那些被弄得气馁或受了羞辱者的数学“创造力”产生多大影响,反而成了全然次要、甚至不足挂齿的问题。

“拒受遗产——或矛盾的代价”那条注,是我头一份写下的反思:盘点经年累月、此处彼处以断片形式回到我这里的东西——既关于“行内现状”,也关于那位我认识得那么好又那么少的人的著作。那也是头一回,我在一眼之间终于看清:那份他携行了大概十五年以上的拒绝,在其数学著作本身中付出的全部“\textbf{代价}”、造成的全部重负。写那条注时,我其实还“落在后面”:因为两年以来(而且没有“人”认为值得通知我),Motif早已走出它们被囚禁十二年的秘密……而今天,在我写下这场对我数学家过去的反思的这一(我相信)最终阶段之时——在我\dmark{}\origpage{230}大致了解那部标志着这次悄然“重返”的值得纪念的文集之后两天——对这份沉重负担的知觉,变得触目惊心:那正是生来该飞翔的人所乐于日复一日、千回百转拖着的重量——一次轻柔轻盈的飞翔,欢快而无畏,迎向未知,为了他的喜悦,也为了托起他的风的喜悦……\fnote{21}{(*)}{我绝无意暗示:“飞翔”、发现世界,是少数特殊者的特权;肯定,我们人人天生受此召唤!只是这份能力,很少有机会哪怕在极有限的方向上(比如数学工作)绽放。但在某个人身上,我曾有幸看到这样一份能力(在“数学”的方向上)格外夺目,如奇迹般被保存,其后却随岁月而倒退。}

倘若他不飞翔,倘若他满足于做一个受人仰慕与畏惧的人,靠不断积累自己优于他人的证据,我无须为他担忧。倘若他拖着他乐于拖的重量,他肯定从中找到满足——一如我自己曾乐于拖重量,并且今天仍在拖那些我在路上尚未学会割舍的重量。我所能给他的,无论好的一面还是坏的一面,他取其所愿。他的选择只属于他自己,我无须为其担忧;甚至无须在此裁决它们是最好的还是最坏的(注62)。对一个人是“最好”的,对另一个人是“最坏”——有时对同一个人也是(只要他肯变,这诚然少见……)。

但我们所做的选择、表达这些选择的行为(即便我们的言语常常否认它们),我们都须自行承担风险与后果。它们常带给我们所期待的报偿(我们把它当作“最好”来领受);而这些报偿本身,有时终于显出其反面(我们把它驳回为“最坏”,甚至当作一桩冒犯)。当人们终于懂得:反面并非冒犯,往往便把它看作一笔该付的价钱,付得老大不情愿。但也会有这样的领悟:某些反面,并不是无情的收银员——为了快活时光,不情愿也罢,总得付账;它们是耐心而执拗的信使,不倦地回来,永远带着同一个讯息;一个不合时宜、不断被拒的讯息——因为比反面本身更甚的,是它那份谦卑的讯息总遭驳回,那份讯息在我们眼前才显为“最坏”:比一千个反面更坏,常常比一千次死亡、比整个宇宙的毁灭更坏——对整个宇宙,我们已经毫不在乎……

\dmark{}\origpage{231}终于有一天,当我们乐于接纳那份讯息,眼睛霍然睁开,看见:曾被畏为“最坏”的,是一场\textbf{解脱},一场宏大的释负——而那份压得我们喘不过气、我们蓦然卸下的重量,恰恰是昨天我们攀住不放、当作“最好”的东西。


\subsection{事件}

\notehead{62}{}

(4月21日)人们会对我说:既然我无须为他担忧,那为什么还要为一个只关乎当事人与我俩的个人关系,铺陈上这么多页!

我之所以感到需要这场对一段关系的某些重要侧面的回顾性反思,是一个具体的、与我切身相关的事件使然(尽管我是带着两年的时差才得知它的)。另一方面,这一事件处在公共领域之中,比那些成名的数学家(如Deligne,或我自己)对待名气较小者或初学者的日常行为举止,还要显豁得多(虽然后者对他人生活的影响,往往与眼下这个个案完全不在一个量级)。所说的事件(即Lecture~Notes的LN~900那部“值得纪念的文集”的出版,别名“下葬文集”),连同围绕着它的一切,在我看来——无论对错——是\textbf{不健全的}。在我看来,给某些来龙去脉作一份周详的见证,直探我今日所感知的事情的根底,对所有人都是有益的,首先要从“当事人”本人做起。经由这一见证与这场反思,我并不想说服任何人相信任何事情(那实在太累人,何况毫无希望!\fnote{22}{(*)}{(5月25日)如果我在这里觉得有必要重申,想说服人是“太累人”而且“毫无希望”的,那恐怕是因为:在我身上某个地方,说服的意图毕竟实实在在地在场,而且也被觉察到了。从4月19日(我得知“值得纪念的文集”LN~900那一天)到4月30日的整个反思,都带着一种内在紧张、乃至分裂的状态——面对一桩全然出乎意料的“事件”的冲击,我尽力消化它带来的讯息。这一紧张,最终随着4月30日的“事物的回归”一注(n\textsuperscript{o}~73)而解开:反思终于回到我自己的身上,随即为我提供了这则讯息的明摆着的钥匙。}),而只是想去理解那些我身陷其中的事件与境况。倘若它们能促使别人在老生常谈之外作一番真正的反思,这份见证就不算白写。

\subsection{驱逐}

\notehead{63}{}

\pmark\origpage{232} (4月22日)这篇文章\fnote{23}{(*)}{即Deligne论谱序列退化与Lefschetz定理的文章(Publications~Mathématiques~35, 1968),“罐装的权与十二年的秘密”一注(n\textsuperscript{o}~49)中已提及。}1968年发表于《Publications Mathématiques》,即我离开数学家世界的前两年。它的出发点曾是一条猜想——我同Deligne谈起过:关于谱序列的一个退化性质,在当时看上去相当不可思议,却经由“算术”的途径、作为Weil猜想的一个推论,变得可信起来。这一动机本身就有极大的趣味,因为它表明:从隐含在Weil猜想里的那部“权的瑜伽”(这部瑜伽最早由Serre在某些重要方面隐约瞥见)里,人们能榨出多少东西。从那时起,我就经常把它应用于各种类似的情形,从“算术的”论证榨取“几何的”(对代数簇的上同调而言)结论。只要Weil猜想尚未确立,这些论证就始终是启发性的,但它们仍然具有巨大的证明力,代表着一种头等的\textbf{发现手段}。Deligne对这条特殊猜想的“几何”证明——借助Lefschetz定理(当时只在特征零确立)——在一个全然不同的方向上另有其趣味,且不提它的头一项功劳:不依赖任何猜想。两条进路所指示的那种联系——一边是Weil猜想(以及当时在我看来其最迷人方面的权的瑜伽),另一边是Lefschetz定理——联系着两件看似互不相干的事情;这联系本身就极富教益。

对我眼下的处境而言,这里有趣的一点——直到今天它才向我显出全部意味——是:这篇文章的读者,几乎无从猜到我参与了其主结果的最初动机,更绝无可能从这篇文章里得知那动机是什么。(另可参见注49开头。)为这样一个结果作陈述,\textbf{自然而然}的路数(我确信,连作者本人最初也是这样)本应是:从那条(的确引人注目的)猜想出发,指出找到的第一条理由——它同样引人注目,而这正是终于把那部著名的权的瑜伽“推销”出去的大好机会,它本身的分量远比该文的主结果更重\dmark{}\origpage{233}\fnote{24}{(*)}{恰恰是这部瑜伽,此后六年里(据我看来)一直保持着秘密!(6月7日)而且(如后来所显现的),它当时被Deligne“记在自己名下”地呈献出来,对Serre与我都不置一词。(见注n\textsuperscript{o}~78\textsubscript{1}、78\textsubscript{2}。)};然后接着谈“Lefschetz定理”的观点\fnote{25}{(**)}{(6月17日)用(“奶牛”)Lefschetz定理来证明谱序列的退化,这一想法出于Blanchard;不过他只在一条苛刻得罕见地得到满足的假设之下——即纤维的有理上同调构成的局部系统是平凡的——才得到退化定理。我熟悉Blanchard的工作,自然向Deligne谈起过;他因此从Blanchard的想法取得启发来完成他的证明,尽管他没有读过那篇文章。Serre对Blanchard证明的记忆比我好,他向Deligne指出:他的证明实际上是Blanchard证明的一个不费力的改写。Deligne在其评注2.10里承认了这一点。但这条引用了Serre的评注,写法却给人这样的印象:仿佛他是事后才得知Blanchard的想法的——事实绝非如此。于是,这篇文章的两条主要\textbf{源头}都被藏掖了:一边是\textbf{算术的动机},它使人得以预见Blanchard结果的一个重大强化;另一边是Blanchard的\textbf{证明想法}——他把后者优雅地加以改写,得到一个Blanchard恐怕不敢指望的结果,也因此从未试图用自己的方法去“得到”它。},它使人得以在略更一般的条件下(基底概形任意,不必真且光滑于一个域)证明最初的猜想,但只限于特征零。而文章实际采取的陈述,却从同调代数的一般的漂亮道理开始(想必十分漂亮,而且以作者惯有的优雅陈述出来)——那些道理,他想必早已同所有人一样忘光了,俨然是Lefschetz定理公理化的路数。主结果(当然,也是唯一人人记得的结果)以系X的面目出现在文章中部;而快到结尾的某处,在“评注2.9”里,读者不知为何,“权”这个词与我的名字才被提起……

文章发表时给我的印象,我已毫无记忆——我在局中,想必只是匆匆浏览了一眼。我肯定感到了一种“保持距离”的意图,但也感到:我的朋友不愿冒被人看作某位“大师”的信徒(或“驹子”)的风险,这原是再自然不过的事\fnote{26}{(***)}{(5月26日)关于我身上的这种态度,可参见紧随本注之后的“高升”一注(n\textsuperscript{o}~63$'$)。(6月8日)把它与我朋友那套据有他人观念的独门手法放在一起看——这里就是第一个典型的例子——我进而明白:我朋友的动机,根本不是要在一个显赫的“大师”面前保持什么“自主”,而恰恰是要藏掖他人的观念在其观念之生成中的作用,以待来日把这些他人的观念也一并据为己有(在第二步里)。(对此可参见“变戏法的人”与“据为己有与轻蔑”两注,n\textsuperscript{o}~75$''$与59$'$。)至于在我的朋友身上,这一癖好得以畅行无阻地发展,我自身应负的那一份责任,可参见“高升”“暧昧”与“格格不入者”三注(n\textsuperscript{o}~63$'$、63$''$、67$'$),其中显出我对这位光彩夺目的年轻人Deligne所曾表现出的某种纵容。}。诚然,倘若他心中有那份对自身力量的安然自信\dmark{}\origpage{234},他原会毫不犹豫地写出一部分量更重、对人人更有用(肯定也对他自己更有用)的著作,而不必担心别人看不出他的本色……(注65)

次年,他第一部大面积的著作——论混合Hodge理论——发表时,情形也多少类似。(我当时认为,这项工作的分量,可与Hodge理论本身相比,把它看作一套“Hodge--Deligne系数”理论的起点——那套理论可惜始终不曾问世……)如我所说,对他和对我都再明显不过:这项工作的“动机”在于我此前数年抵达的Motif瑜伽——它是使这部瑜伽获得有形实现的第一番尝试。在他的著作里点明这层联系,依我看(当时想必也这样觉得),本会立刻使这项工作获得比它凭自身优点已具有的更宽广的分量;同时,这又是一次机会,把读者的注意力引向Motif的实在——它在Hodge结构的实在之后,随处可以感觉到(注63$_{1}$)。

只有带着事后的认识,这些省略才在六年的权的瑜伽沉默\fnoteplainmark{27}{(*)}{(1985年4月19日)关于“六年”与“十二年”的更正,对权而言见第302页脚注(***)(1985年4月18日所写的部分)。}、十二年的Motif沉默(不妨说,是禁令)\fnoteplainmark{28}{(*)}{(1985年4月19日)关于“六年”与“十二年”的更正,对Motif而言见子注“预先出土”(n\textsuperscript{o}~168$_{1}$)。}之后,在这场以LN~900下葬文集为形式的罕见的“重返”,在Hodge--Deligne理论惊艳的开局之后的停滞之中,显出全部意味……但没有人能怀着殡仪人的心肠做出大事!

无论如何,倘若1970年我离开IHES时心智更成熟一些,那么从那一刻起,我就会一清二楚:在那个五年间一直扮演着我最亲近之友的角色的人身上,存在着一种对我而言深刻的暧昧。况且,在一个平静安适的机构里,在同侪间彬彬有礼的情面外观之下,我的离去,最终对人人都是方便的——缘由我自信事后已能辨出,而且各人并不相同。\dmark{}\origpage{235}显然,这一离去对我的年轻朋友再合适不过:他新近就了位,本只需同我站在一边(面对另外三位终身同事迟疑的冷淡),就足以翻转一个未定的局面。当时我不明白这一切意味着什么,是因为我分明不愿去理解那些其实足够清楚、甚至足以自明的事情!一如我生命中常常发生的那样,那时我心中有一种焦虑(我从不曾以这个名字称呼它!),它向我提示:在一个有形而简单的实在,与一幅我不肯放手的实在图景之间,出现了“脱节”:即我在即将告别的机构里所扮演角色的图景,而更甚的,或许是我与朋友之间那层关系的图景。正是这种拒绝面对一项无可置疑的实在,以及作为这桩矛盾之征兆、我却死死攀住的焦虑,使得那次“有益的挣脱”的情节,在当下显得如此痛苦\fnote{29}{(*)}{关于这一情节,见注n\textsuperscript{o}~42。}。

老实说,由于从未曾为这段关系专门写过一篇反思(除了一些散见于给我朋友的书信里的反思的开端,而且没有一封得到回音……),我此前不曾注意到:我的朋友对我之关系中矛盾态度的最初征兆(诚然隐约,却绝不会错)至少可以上溯到1968年,即“大转折”前两年。那是一个这段关系显得完美无缺的时刻:在数学层面上是无云的契合,处身于一段单纯而亲切的友情之中。人们此刻尽可以讥讽那些关于天真、创造的孩子的美好“长篇大论”!

然而我深知,那种契合是一个\textbf{实在},绝非幻象;一如那“精微的东西”是实在——那股创造的力,其后的作品只给了它一个苍白的反照。“天真”与“冲突”是两个可以触摸的实在,略一警醒的知觉就能辨认,绝不是什么概念;而且在我看来,它们按其本性相互排斥,一个排斥另一个。然而毫无疑问,这两个实在,在不同的层面上\fnote{30}{(**)}{在另外两三次机缘里,我得以确认:在同一个人身上、在某个时刻,竟有这等共存——包括我自己的身上,在有些时刻。},共存于我朋友对我的关系之中。在我所说的那个年代,“冲突”似乎并未干扰数学的创造性——至少在\dmark{}\origpage{236}孤独中做的工作里、或在面对面的交谈中进行的工作里没有。诚然,在我方才谈到的那两篇文章里——它们毕竟是这项工作最切实的果实之列——“冲突”的印记已经分明可见。而带着十五年的事后的认识,经由这些天与这些星期的反思,我看清了:这道印记(无论多么隐约),以一种触目惊心的方式,预演了冲突将逐步压倒那股最初的冲劲、并在岁月里剥除其最稀有的本质——那造就大命运的东西——时所采取的特殊形式\fnote{31}{(*)}{这一节如此高贵的抒情,使我有点脱离了尘世的实在。倘若我在这里称这道“印记”为“隐约的”,那是因为我自己还裹在一层厚茧里,还难舍那副我视若珍宝的眼罩!一旦终于摆脱了它,我便看清:所说的“印记”,是一场粗劣的藏掖;我不肯看见它,是出于我身上的某种纵容——6月1日的“暧昧”一注(n\textsuperscript{o}~63$''$)里,我对此有清楚的认识。至于“冲突对我年轻而光彩的朋友那股最初冲劲的压制”,我几乎把它说成一场令人惋惜的宿命,仿佛这位可怜人是全然不自愿的受害者,从而——唉!——也一并取消了“大命运”的收益。然而,他对自己的命运,负有与我对我自己的命运同样的责任。倘若早在我离去之前,他就选择了做他导师的掘墓人(权且从这里开始)的角色,倘若种种情势(其中包括时代精神)有利于这一选择,把“打谁都行的大老板”的角色让他尽情享受,那么,他也同样选择了把威望与权力所能给人的特权——包括(悄悄地)碾压与巧取的特权——品尝到最后一滴。鱼与熊掌不可兼得;按照事情的本性,经由这一选择(在这选择里,他自有一些好同伴),他失去的是那些更精微、更不趋时的东西的收益……(6月初的一则未注日期的脚注)。}

\notehead{63$_{1}$}{}

(5月26日)亦可与注60末尾脚注(*)处的评语对读:那里确认了Hodge--Deligne理论的自然的发育所遭到的“封锁”——由于对我引入的某些观念核(这里是六种运算——Motif与它们密不可分)的拒斥态度;那态度与这里考察的同属一类,因此早在“Hodge理论I与II”发表时就已显形。

同样这种态度——竭尽所能(乃至过之!)抹去我的影响的任何痕迹——也见于与Mumford合作的那项工作(注n\textsuperscript{o}~47已提及):论模多重簇的Mumford--Deligne紧化。(这项工作同样先于我的离去。)它利用一条从$\mathbb{C}$上的拓扑结果(经由超越途径得知)过渡到特征$p>0$的结果的原理——那条原理,是我在五十年代末为基本群理论引入的\dmark{}\origpage{237}。早在六十年代初,我就建议利用这一方法,在全特征下证明模簇的连通性\fnote{32}{(*)}{(1984年9月)核实之后,这一情节确确实实在前引工作的导言里(第75页)得到了记载。}。然而这一想法碰到了一些技术困难,它们曾使Mumford却步,而在他们那项工作里被优雅地克服:办法是引入\textbf{模多重簇},以及一种具有完美性质的“紧化”。模多重簇这个想法本身,至少“在字里行间”,可以在我在Cartan研讨班上的“Teichmüller”报告里找到——那还是site与topos的语言尚不存在的年代。Deligne使用的语言本身——在他本可使用一整套为表达这类情形量身定做的site、topos、多重簇的语言的地方,却用了“algebraic stack”——相当清楚地显示了(带着事后的认识,并且借着后来规模大得多的几场“战役”之光照亮)抹去这项光辉工作中所运用的某些主要观念之来源的意图。肯定正是这种态度(我在“拒受遗产——或矛盾的代价”一注,即n\textsuperscript{o}~47里,头一回预感到它)产生了那记“电锯效应”,斩断了其后对模多重簇的反思——尽管在我看来,模多重簇属于迄今所提炼出的一切“具体的”数学对象之中最美、最基本之列。

顺带指出:我在五十年代末引入的那些论证,使人得以(有赖于Mumford--Deligne紧化)不仅在全特征下证明模多重簇的连通性,而且确定它们的“素于$p$的基本群”:即寻常Teichmüller群的“素于$p$的投射有限紧化”。

\subsection{高升}

\notehead{63$'$}{}

(5月10日)带着再多不到三个星期的事后认识,我如今看清:我当时那种自以为“体谅”的态度——对着那“再自然不过”的保持距离的意图——实际上是一种洞察力的欠缺,是对我年轻而光彩的朋友的一份纵容。倘若那时我信任自己健全的知觉,而不肯被那些含糊的老生常谈弄得眼花缭乱、哄骗自己——它们摆出“体谅”的姿态,甚至摆出“宽宏大量”的姿态(“总不能因为他不把我的名字捧出来,我就说他几句……”)——我那时就会看清我如今(十六年后)看清的东西。我可以称之为:对读者、对我自己、也对他自己,缺乏一份正直。简简单单地看事情,毫不害怕把它们\dmark{}\origpage{238}叫出它们的名字,我本会有能力像现在这样坦然谈论它们;而我的朋友当时本可从中取法——或至少他会懂得:即便凭他那般手段,他的长辈们(或至少其中一位)也期待他在工作中具有与他们自己相同的那份正直。于是我看清:在那次机缘里——它在我退出数学舞台之前,彼时我绝非“局外”,而且对年轻朋友无疑还有某种道德上的分量——我经由当时表现出的这份\textbf{松懈}\fnote{33}{(*)}{(5月28日)这里用“纵容”一词,比那个有点躲闪的“松懈”,更能表达我那份态度的性质。在我与我年轻而光彩的朋友的关系中的这份纵容,在昨天的反思里显现得更清楚;见“格格不入者”一注,n\textsuperscript{o}~67$'$。},没有尽到我对他的责任。这一松懈,在“Hodge理论II”发表时得到了印证:那是Deligne的学位论文工作,其中对Motif、对我,都不置一词。诚然,到那时,数学和我朋友其人,对我都已十分遥远,仿佛隔着一层浓雾!

从我得以观察到的我朋友的演变——精神上的与数学上的(这两个侧面血肉相连)——我如今看清:在我与他相识、并被他的智力手段、他的目光的敏锐、他在数学上理解的敏捷所打动的那个时刻,我丝毫没有在他身上辨出成熟的欠缺;也没有(后来)辨出,他社会地位的扶摇直上——在不过四年的工夫里,从一个无名的学生,一跃而成为数学世界的明星、终身教授,在一个已然声名显赫的机构里被授予可观的特权与权力——将对他产生怎样的影响。我不后悔为他铺平了这条路,还让它走得更快——但我确认:那是出于我自身的眼力与成熟之不足。我为他尽的这份“效劳”,算不得效劳。只要我的朋友自己尚未把那场收成——他在我漫不经心的协助下为自己预备的收成——进行到底,这就不曾是一场“效劳”。

\subsection{暧昧}

\notehead{63$''$}{}

(6月1日)自从在我与我的朋友皮埃尔的关系中确认了那份“松懈”(或按其间出现的更贴切的说法,那份“纵容”)以来的三个星期里,我已有机会在反思中更清楚地认识到我身上的某种不够严格、某种纵容。它们最初显现于我与那个人的关系——在我那里,我比对任何人都更把他当作“格格不入者”来对待——但也显现于我与另一些数学家的关系,在后者眼里我扮演着长辈的角色。迄今我在这个方向上侦察到的\dmark{}\origpage{239},表现为我身上的一种暧昧——恐怕也在那位扮演学生角色的人身上:在那些场合,他把得自我的观念与方法、乃至我交给他的一项完整工作的详细施工图据为己有,而不明确指出来源,有时甚至只字不提。这类情形相当常见,既见于六十年代,也见于我离去之后直到最近这几年的年月。我觉得,在所有这些情形里,在某个层面上,我感到了那份暧昧;它表现为一丝不安的阴影,直到最近这些天才头一回被检视。使我入局、参与某种心照不宣,并使我掠过那丝不安、从不予理会的动机,在于\textbf{顺应}我对自己、对所谓“宽宏大量”应有的样子所抱的某种形象。真正的宽宏,并非生于一种顺应成见的心态,也不是出于想要成为“宽宏的人”、并在自己与他人面前显得宽宏的心思。被压抑的不安,每一次都是明明白白的征兆:这份“宽宏”是做作的,是一种\textbf{姿态},而不是宽宏之真正的、毫无保留的自发赠予。

在这份不安里,我辨出两种来源不同的成分。一种来自“老板”、来自那个始终不得满足的“我”:他没能两头都赢——既分享一项他知道其中有自己(或大或小)一份的工作的信誉,同时又配得上某种品牌形象——其中赫然印着“宽宏大量”这条老生常谈的标签(此外还有许多别的标签)。另一种成分来自“孩子”、来自我身上那个不为姿态与门面所欺的人:他朴素地感到这境况的虚伪\fnote{34}{(*)}{(6月5日)当我在这里说,不安(部分地)来自“孩子”时,这种说法给人的是一幅不真实的图景。并不是对一个虚伪境况的天真知觉造成了什么不安。不安是一个\textbf{抵抗}的征兆——抵抗那份知觉,是(在某些层面上分明被知觉到的)实在(这里是:一个虚伪的境况)与我所攀住的实在\textbf{图景}之间的“脱节”(在这个案子里,图景即:我正在“宽宏大量”,而且不能不这样!)——我为了图景的利益,把那份不合时宜的知觉\textbf{推开}、压抑。在这个案子里,一旦我放弃抵抗、容许知觉进入意识之眼的视野,“不安”便随之止息,连同那个虚伪的境况。我本想说:“这是就涉及我当下的虚伪境况而言,而不是就处在过去的境况而言。”但转念一想,我发现:我方才谈到的那些“过去的”虚伪境况,直到今天仍以本来的面目在身,或者至少直到三天前的反思——仅仅因为它们从未被检视、从而从未被解决。我成了它们的囚徒,以至于一有机会,就机械地重演同样的境况。我对自己的“观”的力量的知识(我在“欲与观”一节,n\textsuperscript{o}~36里谈过它)当时对我毫无用处,因为我缺乏对日复一日身陷其中的境况、对知觉与其“拣择”的不停息的游戏的留意——那孩子与老板的游戏使知觉沉默……}。不但对我自己虚伪,而且\dmark{}\origpage{240}对他人也虚伪。

总之,我的“宽宏”在于入局:入一场游戏,游戏中对方把他得自他人的观念呈献为自己的,也就是呈献一幅他自己与某种实在的图景——而他和我都心知肚明,那是假的。于是我们成了同谋,共谋一桩可称之为“作弊”的事,双方各得其所。至少按“我那个年代”通行的共识,这是一桩“作弊”——而那些共识,依我看,今天人们仍在嘴边念叨。倘若被当作“我庇护的人”的发现来使用的是别人的观念,我肯定不会入这样的局\fnote{35}{(*)}{“我庇护的人”这个说法,曾被我昔日的一个学生用来称呼我当时的、刚在数学里做出漂亮东西的一个学生;它曾使我咬牙切齿。然而,细想之下,我此刻正在考察的这份暧昧的境况,确立的恰恰是一种虚伪的关系:其中一方俨然是另一方的“受庇护者”。}。然而,我给予默许,让诞生于我的观念被呈献为他人之物,并不改变事情的本质——唯一的区别是:这回我们是两个人在作弊,而不是一个人。撇开涉及我个人的这一面(我自己参与一场作弊,一种违背我自己也自认服膺的共识的行为)不谈,有一点再清楚不过:怂恿他人作弊(即便这作弊看似只由我们自掏腰包——而事实绝非如此),或至少怂恿他人对一个他自己也假装服膺的共识持一种暧昧的态度、实则阳奉阴违,这里没有任何宽宏可言。真正的宽宏,本性上造福所有人,首先要造福它于其中显现的那个人、它所施与的那个人。我的这份暧昧态度——在人身上诱发或怂恿暧昧,并使我得以摆出“宽宏”的架势,而按正常的道理,对方这时应当显出几分作弊者的模样(而事实上我们俩都在作弊)——这份态度,对我既非善举,对他亦然。

只需把事情检视一番,明摆着的道理就会显现,甚至无须求助于经验、求助于“事件的教训”。然而,终究是事件把我引到了这番检视,使我终于发现一个明摆着的道理——三十\dmark{}\origpage{241}年前我同样有能力发现它,那时地平线上还没有出现一个学生,要跟我学一门手艺,并在与我的接触中浸染于这门手艺练习中的某种精神。我已有机会谈到工作之中的那种“严格”——我相信自己做出过表率(见“严与苛”一节,n\textsuperscript{o}~26)。但今天我还察觉到,在“工作”本身之外,也有一份严格之不足,表现为暧昧,表现为我所说的纵容。我觉得,我的这份暧昧,不是从我的任何一位长辈那里传来的——他们(我相信)对我都有与自己对自己相当的要求。在这一特定态度的暧昧之外,我还在我这个人本身之中辨出一种暧昧——在《收获与播种》第一部分,我已有不止一次机会谈到它。这份暧昧,随着1976年禅观的发现开始消解;而它的某些征兆——表现为已然成习的态度与行为(尤其在我与学生的关系里)——想必一直延续到今天。

显然,我身上这份暧昧,在我的某些学生那里找到了肥沃的土壤。那种经由默契做成的事,如今似乎已成了“上流”数学界习俗中的一个基调:在那里,浑水摸鱼(有或没有“当事人”的同意)、乃至有组织的劫掠(当放胆为之的人属于那不可撼动的精英时),似乎已成了寻常得无人再表惊讶的惯例——而同时,人人也都绝口不提。我身上的“老板”很想划清界限、揭发、愤慨——可这样一来,我不过是在延续我身上那份同样的暧昧,它的丰饶收成,我今天有目共睹。

\subsection{同伙}

\notehead{63$'''$}{}

(4月24日)\fnote{36}{(*)}{本注出自“本能与时尚——或强者的法则”一注(n\textsuperscript{o}~48)的一条脚注——在那则注里,我曾断言Verdier关于导出范畴的工作从未发表,却不知道其学位论文的一个“第0状态”已于1977年发表。至于Verdier围绕着那本应构成其学位论文工作的理论所玩出的那些古怪的腾挪,其全貌可参见“赊账的学位论文与全险”一注,n\textsuperscript{o}~81。}两天前翻阅刚收到的Mebkhout的一份抽印本,我撞见一处对 J.L.~Verdier一项工作的引证,题为“导出范畴,第0状态”,发表于SGA~4½(Lecture~Notes n\textsuperscript{o}~569,第262--311页)。我到今天才得知这一发表,情有可原:此前我从未有过捧起这一卷的荣幸——无论Verdier,还是其作者Deligne\dmark{}\origpage{242},都不曾费心在出版时或后来寄我一本。我不知道,当年与我共同组成评审委员会、仅凭一篇17页的导言(至今未发表)就把“理学博士”头衔授予 J.L.~Verdier的 C.~Chevalley与 R.~Godement,十年之后,有没有份收到这部不同寻常的“学位论文”的“第0状态”(这一回有50页)!我记得有一天曾亲手捧读过一部几百页的严肃奠基工作,它满可以理直气壮地算一篇优秀的博士论文,并且大体对应于我1960年前后向Verdier提议的那项奠基工作——只除了一点:彼时“三角范畴”框架(他用来表达导出范畴的内蕴结构)业已显出不够用。

几乎无须说:我的名字,在这份学位论文的“第0状态”里无处可见。人们纳闷:它到那里去做什么。众所周知,导出范畴是Verdier引入的,为的是让他得以发展拓扑空间的所谓“Poincaré--Verdier对偶”与解析空间的所谓“Serre--Verdier对偶”;只待某位当值无名之辈\fnote{37}{(*)}{关于这位可疑人物的几点情况,可参见“当值无名之辈与好上帝定理”一注(n\textsuperscript{o}~48$'$)。}出头,为他自己发展一个把二者综合起来的理论,并且理所应当地(无名的学生至少也得做到这一步!)被称作“Poincaré--Serre--Verdier对偶”。经过这一切,我剩下的只有随大流,做几处势在必行的调整,以便在概形的平展上同调或凝聚上同调那个再特殊不过的框架里,发展Poincaré--Verdier对偶与Serre--Verdier对偶……

我刚刚才得知(图书馆到底有用!)SGA~4½的内容\fnote{38}{(**)}{关于这一卷,可参见“白板”一注,n\textsuperscript{o}~67。}——在那里,人们还不吝让我以合著者、或不如说Deligne的“合作者”(原文如此)的身份露面(既不费心告知我,更不费心征求我意见)。显然,这是五年后问世的“下葬文集”的先声——那部文集,几天前我有幸拜读(见由这一事件激发的注n\textsuperscript{o}~50、51及以下)。但我无需亲手捧起那部“预葬”文集,连同那篇不肯自报家门的幽灵论文的铁证,便早在去年就已懂得:这份“学位论文”的下一个状态,除了我,不会有别人来写。于是我就着手这项工程,即\textbf{《追寻栈》}\dmark{}\origpage{243},从我那位显赫的前学生十七年前停下来的地方接着做。

\subsection{授职}

\notehead{64}{}

(4月25日)然而昨天,我在系里的办公室里找回了一份副本。其实是前后相隔一年、写于1968年4月(?)与1969年4月的两份报告。我用十七页的篇幅,评阅了他在IHES三年科学活动中进行的十五项工作。其中有关于Ramanujan猜想的工作(原书拼作Ramanuyam),关于模site的紧化的工作,以及Hodge理论的推广。报告里评阅的这些著作(单是我刚点名的几项就足够了)见证了一种惊人的创造性,以完美的从容、仿佛游戏般地铺展。把Weil猜想的证明——那还是紧随这第一次向未知的冲刺之后——放到一边,我觉得,他后来的著作,只给了那次独一无二的起飞一副苍白的图像:一个手段出众的年轻心灵,享受着使其绽放的种种同样出众的条件。然而,不得不相信,这些“出众的条件”中的某种东西,必定喂养了另一股力量——一股与求知的冲动异质的力量;它最终反客为主、取而代之,使最初的冲劲改道并被吞并。而显然,这个“某种东西”与我这个人有关……\fnote{39}{(*)}{(5月26日)至于我身上的某种纵容——它喂养了这个“某种东西”——可参见比本注晚两周的“高升”一注(n\textsuperscript{o}~63$'$)。}

这篇简短的评述报告(我打算作为附录收入本卷)从不止一个角度看都饶有趣味,包括数学的角度(评阅的有些著作,至今仍未发表)。在报告的几处,我预言:Deligne只勾勒了轮廓、处理了关键之点的某些工作,将由未来的学生们去发展。这些学生从未出现,缘由是他后来对常人世界的关系起了变化\fnote{40}{(**)}{在我还常在IHES与他为伴的年月里(尤其在我的研讨班上),Deligne对其他数学家、尤其对来参加研讨班的青年研究者(往往是初学者)的关系,带着亲切。我在那里看到的对他人思想的开放——哪怕它表达得笨拙甚至混乱——同我们数学的两两对谈里一样。他有一种能力:循着他人的思想,进入他人的图像与语言;这能力我一直欠缺,而且(依我看)使他比我远更适合“老师”的角色——善于激发他人身上天职与创造性的绽放。}。在我评阅的那些想法之中,据我所知\dmark{}\origpage{244}唯一曾被别人发展出来的(那人由此俨然成了Deligne的学生),是上同调下降理论——由圣多纳在SGA~4里发展(因此还处在最初冲劲的时期),此后成了上同调武库里最常用的工具之一。有一处既有趣又有代表性的细节:后来成为Deligne文章的四项工作里有三项\fnote{41}{(*)}{四项工作中唯一不受我直接影响的一项,是从Weil猜想推出Ramanujan猜想的工作;它处在一个研究方向(模形式的方向)上,那是我数学修养里最严重的“空洞”之一。其余三项是:Leray谱序列的退化,Hodge--Deligne理论,模多重簇(与Mumford合作)——见“驱逐”一注(n\textsuperscript{o}~63)与子注n\textsuperscript{o}~63$_{1}$。},我在行文之间,以动人的细心,让读者感觉到这些工作与我引入的想法、我提出的问题的关系——仿佛要抢在前面,预先对付作者日后在他的文章里对这些内容的沉默(这几篇文章,在我写报告的时候,没有一篇已发表,而且我相信,甚至还没有成文)。

\subsection{死结}

\notehead{65}{}

(4月26日)同样清楚的是:把一部大面积的“瑜伽”(权的瑜伽,进而Motif的瑜伽)留在自己怀里——关于它,我对除他之外的人也不止一次谈起过,但唯有他亲密地消化了它、把握了它的全部分量——这本身授予他一种额外的“优势”:作为一件无可比拟的发现工具的独家持有者,用以理解代数簇的上同调。然而我不认为,在我还完完全全在场、活跃于数学世界、毫无迹象预示我无限期离去的那个时刻,这一诱惑起过什么决定性作用。它必定是随着、或在我的离去之后才出现;我的离去,是攫取一份遗产的“天赐良机”——那份遗产本是完全理当归他所有!——只要藏起遗产,连同它的来源。

正是在这里,我又一次看见那个深藏的矛盾的死结,在一个极端而格外触目的个案里,远远超出一切个案:我想谈的是无知、不屑、深深埋藏的怀疑——它们围绕着我们自己身上所怀的那股创造之力:那笔独一无二的遗产,比一个人所能传递的一切都更珍贵。正是这种无知、这种对我们身上最珍贵、最稀有者的阴险的异化\dmark{}\origpage{245},使我们得以嫉妒在他者身上知觉到的力量,使我们得以贪求那力量在他人身上的果实与外部印记——那股我们已在自己身上遗忘的力量。一旦这份嫉妒、这份\textbf{取而代之}的欲望扎下根来、找到滋蔓的机会,把本可用于创造性绽放的能量引向这一欲望,我们身上的异化就变得更深,安营扎寨。我们越接近“取代、排挤、炫目”这个觊觎的“目标”,就越远离、越割断我们身上那股精微的力量,也就剪断我们自己创造冲劲的翅膀。在我们顽强地抬高自己的持久努力中,我们久已忘记飞翔,忘记我们生来就是为飞翔的。

在与我的关系里,从我们相识那天起,我就感到我的朋友完全自在,没有任何征兆使我疑心:他会为我的名声或我这个人感到丝毫的震撼或目眩,或者他心中有任何未曾说出的怀疑——无论是关于他在数学领域的天赋,还是关于任何别的。同样,依我看,他也确实从我和我所处的环境——包括我的家庭——得到了一种友好的、亲切的接待,足以使他自在。但那份吸引我、也吸引别人的单纯而似乎毫无疙瘩的自然,肯定不是等到这次相识才出现并绽放的。他这个人散发的、使他如此可亲的印象,是一种和谐的平衡:他对数学的爱好,丝毫不带吞噬人心的女神的形象。在他旁边,我多少显得像个不知悔改的“书呆子”,甚至“粗鲁迟钝的家伙”——我还记得,他含蓄地惊讶于我与身边的大自然、与四季的节律之间缺乏深入的接触;而我穿行其间,几乎视而不见……

然而,那份深藏的“怀疑”——当时我绝无能力知觉它(甚至今天,处在类似的境况里,恐怕也未必)——必定在我的朋友身上,早在我们相识之前就已存在。带着事后的认识,我毫无歧义地看到它的第一个征兆早在1968年;其后这些年里,又有更清晰的征兆\fnote{42}{(*)}{(5月10日)实际上,还有一个“十分清晰”的征兆,早在1966年;见注n\textsuperscript{o}~82的脚注(*)(原书第329页)。}。不过这些都是“间接的”征兆——从我第一手观察到的东西里,没有一个以怀疑、缺乏自信的面目出现;恰恰相反,而且年复一年越来越如此:一种自满,一种\dmark{}\origpage{246}蓄意的轻蔑口吻,乃至鄙夷。但这样的“反面”恰恰揭示其对面——它们成双,反面是正面的影子。

我还经由第三者得知:对某位他从未有机会亲近相识的、声名显赫(而且以难相处著称)的数学家,他曾处于巨大的紧张期待之中,生怕那位大人物不把他看作与自己的伟大相称。这一证词,与我在年轻朋友身上亲眼所见的截然相反,以至于我很难相信它(那是1973年)。带着事后的认识,它却与我另行得知的那些分裂的征兆相互印证,而且全都指向同一个方向。这份分裂,以及我在其中扮演的角色——作为一桩在相识之前想必还弥散无形的冲突的某种凝固剂——在通常情形下,在与一个曾(以这样或那样的意义)是“老师”、或至少是传递者或托付者的人的关系演变里,恐怕会一直隐而不显。这样,我的离去,就成了一个所有人都未曾察觉、而我或许是唯一知情的冲突的\textbf{显影剂}。

而我今天的“归来”,则是第二个显影剂,无疑更加不合时宜。除了它迄今已就我自己的过去、我的现在、以及我曾爱过而至今仍与之相连的人们所教给我的,我全然无法想象它还将向我揭示什么。也不知道它将向那个人揭示什么——最近一个星期以来,他正处在我反思的这一最终阶段的中心;上个月我曾给这一阶段命名(而且我没想到说得这样准……),叫作“\textbf{一段过去的重量}”。

\subsection{两个转折}

\notehead{66}{}

(4月25日)我的朋友皮埃尔对我的关系中那种蓄意的轻蔑与对抗,只限于数学与职业的层面。个人的关系直到今天仍是一种情谊与敬重的关系,不止一次表现为打动了我的细致体贴,分明是真实而无保留的感情的征兆。

在我离开IHES之后那些紧张的年月里,IHES渐渐沉入遗忘,一如这一情节长期不为我所理解的教诲。于是又有十多年,我的朋友在我这里(想当然地)依然是我数学上的特权对话者;或者更确切地说,1970到1981年间,他是我唯一的对话者(除一桩情节外):在我数学活动断断续续的时期,每当需要对话者的时候,我想到的就是他\dmark{}\origpage{247}。

也是把他当作离我最近的数学家,在最初几次机缘里(1975至1978年间),我同样自然而然地向他求助,请他为与我一起工作的学生提供帮助、担保或支持。其中第一次,是1975年Sinh~夫人的学位论文答辩——她在越南,在异常艰难的条件下写成了论文。他是我为论文委员会联络的第一个人。他推却了,话里带着意思:这只能是一篇冒牌的论文,他绝无意为之作保。(不过我到底有办法说动Cartan、Schwartz、德尼与齐斯曼的好心,请他们为这出“把戏”助一臂之力——答辩在关切与温暖的同情的气氛中举行了。)在其后三年里,又经过三四次同类经验,我才终于明白:在我那位显赫而有影响的朋友身上,有一种蓄意的对抗姿态,对着我的“1970年以后”的学生,也对着只带有我的影响之印记的工作(至少“1970年以后”着手的那些)。我不知道,我在好几次这样的场合所觉察到的那种公然的轻蔑,是否也多多少少见于他对另一些数学家的关系——那些他视为远在自己之下的人。他引以为荣的那套无所节制的精英主义的精神本身,会使我作肯定的猜想。无论如何,自1978年起,我不再为任何事情向他开口。这并未妨碍他那份令人气馁的力量,又找到机会有效地施展。

大约也是在那一年,头一批征兆——起初很隐约——显现出一种对我自己的数学活动的轻蔑态度。第一次机缘,是我关于胞腔地图的反思,起于一个关于它们的发现,它曾使我如遭雷击(参见《纲领概要》第3节:“与儿童涂鸦相联的数域”)。这项发现(诚然“平凡”,本来毫无使我的显赫朋友动容乃至感兴趣之处)是另一个数学之\textbf{梦}的起点与第一份材料——那个梦,其规模可与Motif之梦相比,只是到三年之后(1981年1--6月)才随《穿越伽罗瓦理论的长征》开始成形。这些笔记与同期其他的笔记(手稿凡两千页)构成对那片“新大陆”的头一次巡视——一桩关于儿童涂鸦的平凡议论,使我隐约瞥见了它。

\dmark{}\origpage{248}在这项紧张的工作过程中,有两三次我给我的朋友写信,把我的一些\textbf{想法}告诉他,有时也提出一些技术性的问题。当他肯对我的问题开口时,他的评点始终同样清晰、同样中肯,见证着他年轻时便已打动过我的那同一副“手段”。但一份自满已经磨钝了那份曾经使我陶醉的求理解的渴望,也磨钝了那种经由“小”事把握大事的能力,那种倾听大小事物而把握或构想宏伟蓝图的能力。那种能力不属于理智、不属于单纯的“效率”、不属于对一门已成学科的“掌控”或对已知技术的“掌控”;它在理智的层面上,映照一件本质全然不同的事——孩子的那份\textbf{惊叹的天赋}。这份天赋在他身上仿佛已经熄灭,仿佛从来不曾有过。至少就他对我的关系而言是如此——而在此之前,对我的“1970年以后”的学生,已然如此。他已成了一个大人物;他对数学的进路,变成了不折不扣的那种“竞技”态度——一两个月前我才头一回考察过它,而且我自己对它也绝非陌路……

或许我本可以给自己找到一个理由,解释那在共同热情中的契合、那条曾经连接我们的深深纽带的显然缺席。我本会很满足——多半如此——在机会来临时,向朋友的巧思、向他关于数学事物世界的广博知识,提出或多或少技术性的问题,或者仅仅请教信息。但就在那一年(1981年),那份轻蔑的征兆突然变得如此粗暴\fnote{43}{(*)}{(5月28日)关于这第二个转折的新的照亮,另可参见“反常”一注,n\textsuperscript{o}~76。},以至于我对再同他交流数学问题——哪怕偶尔——全然失去了兴趣。($\Rightarrow$注67)

\subsection{白板}

\notehead{67}{}

(4月26日)正是在昨天写前面那些文字时,我把我们关系中的这个新转折,与1982年(也就是几乎恰好在这个苛刻的转折关头)Lecture~Notes那部“卓越文集”的出版联系了起来——它为我的数学下葬揭幕,没有鲜花,也没有花冠!既然我已被判定数学地“死去”,我的朋友仍在这里或那里继续回答一些归根结底已无从必要的数学问题,这算得上某种恩典……

\pmark \origpage{249}试着去倾听事件的意味,我感到,同样不是偶然:轻蔑、数学上的漠然(对着那些——他的“健全的数学本能”本该告诉他——灼热而多汁的东西!)的初次显现,至少就他对我的关系而言,差不多就落在“预葬文集”SGA~4½出版前后的时刻,即五年之前\fnote{44}{(*)}{关于此事,可参见前天“同伙”一注(n\textsuperscript{o}~63$'''$)。(6月5日)本注的反思,在本注及随后三则注(“白板”“格格不入者”“绿灯”“颠倒”)里得到继续;它们使人隐约看到“SGA~4½之役”的意味,及其与母研讨班SGA~5之“拆解”的联系。这一反思又在“我的学生们”的行列里、尤其在“我的学生们(1)--(7)”的续篇里重新拾起;在那里,一场对那个我的上同调学家学生们学艺的研讨班的真正屠杀的图景,渐渐显露。此外,整场战役里摊开着一种漫不经心的轻蔑——我对它在同一时期初次显现的那种“隐约的轻蔑”的知觉,在我朋友对我的关系里,只不过是一个非常苍白的反照。还有一桩联想,是一两个星期前来到的,关于我朋友对我的关系中那个“第一个转折”的时刻——1977年底或1978年间。正是在1978年,我的朋友得到了他那枚实至名归的“奖章”(为了Weil猜想的证明)。这一新头衔(系于一桩“困难尽人皆知的”猜想的证明)被我的朋友以何种方式内化,在“葬礼悼词”(关于我这亡故者的文章)及其对偶文章(关于他自己的)里,触目惊心;那两篇文章,诚然是整整五年之后,才在一个“重大场合”问世。对此可参见“葬礼悼词(1)——或恭维”一注,n\textsuperscript{o}~104。}。单是这一卷出版时的种种情势,就已足以见证一种蓄意的轻蔑姿态——既是隐约的,又是炫耀的。不劳烦征求我意见,甚至不告知我一声,就把我列为Deligne的“合作者”,而且小心提防不给我寄一册——单是这一点,在我看来,就比一篇演说更有说服力。更不必说,Deligne这部书,就其主体而言,本是要把我在十五年多以前发展的那些工作——那时我还没有听说过我这位光彩夺目的朋友的名字!——变得让更广大的读者更易亲近。一份轻蔑,以及继之而来的傲慢,必定由两头喂养:一头是我的缺席——它使我对一切毫无察觉,浑然不觉地“照单全收”;另一头是某种气候——它使这类张冠李戴得以“通行”,似乎引不起最轻微的评注。总之,关于这一卷,以及关于他筹备的那部下葬文集,我没有收到任何人(尤其在我自以为在数学家世界里还拥有的一众朋友中间)的只字回音。

\pmark \origpage{250}在导言里,作者更是开门见山地亮明底色。该卷的目的,是让“非专家”免于“求诸SGA~4与SGA~5那些芜杂的报告”,“剪除多余的细节”,“让使用者得以忘掉SGA~5——人们可以把它看作一系列枝蔓,其中一些还非常有趣”(对这些“枝蔓”倒还算客气!)。“SGA~4½的存在,不久将使人得以照原样出版SGA~5”——一句神秘的话,因为人们纳闷:这一出版(出的是一部他们劝人忘掉的东西)拖拖拉拉已有十二个年头,而且拥有一组自洽完美的结果(它们不曾等Deligne来提炼与证明),凭什么要以SGA~4½的存在为前提。在提出这个问题的同时,我也隐约看到一个简单的回答,以及那个可怜的SGA~5研讨班坎坷命运的也许一种解释(注68)(那是我于1965/66年间纵横发展起来的,比DeligneSGA~4½分卷的出版早了十一年\fnote{45}{(*)}{对这一“奥秘”的澄清,可见“绿灯”一注(n\textsuperscript{o}~68)的一条脚注(4月28日所写)。})。当人们被告知(第2页):SGA~5的原始版本里,“Lefschetz--Verdier公式只是猜想式地确立的”(这对Verdier未免刻薄——他本该会证明自己的定理,而那定理比SGA~5更早\fnote{46}{(**)}{(6月10日)对此的更精确说明,见“屠杀”一注(n\textsuperscript{o}~87)的子注n\textsuperscript{o}~87(原书如此,未印附号)。}),而且“此外,诸局部项在那里也未加计算”。这对非专家读者(本卷首先面向的读者)可能显得是一处令人遗憾的缺漏。略知门道的读者则很清楚:这些“局部项”,到今天也仍然没有“计算”出来;而且,倘若有人问那位光彩夺目、武断专行的作者,他在这句话里(就一般情形而言)所说的“计算”究竟何指,他自己也会窘迫万分\fnote{47}{(***)}{(6月10日)在一般的Lefschetz--Verdier公式里,对一个系数层与其自身的上同调对应,那些“局部项”(对应于不动点集的连通分支)在写出公式的一刻便毫无歧义地得到定义。这些局部项的“计算”问题,只在具体的个案里才取得精确的意义;其中最简单的之一,是Frobenius态射的情形——在那里,它们简单地由这些点上的纤维所诱导的自同态的寻常迹给出。这条公式,在口头研讨班里,已作为一个远更一般的公式的特例得到完整的证明。}(不过,显然没有人想到去问这句不客气的问题)。

一句含混的话——“这个研讨班(? )含有\dmark{}\origpage{251}另一个证明,这一份在Frobenius态射的特殊情形下是完整的”——似乎暗示:SGA~5终究没有给出(谁又能料到呢,在一部“枝蔓的文集”里!)它所宣布的那个主要“结果”——一条迹公式,从而蕴含Weil式的L函数有理性——的完整证明;幸好“这个研讨班”前来救场,亡羊补牢,犹未为晚……

在第4页,我们得知,“Arcata”诸报告的目的,是“给出平展上同调诸基本定理的证明,剥去包裹着它们的、SGA~4里肆虐的non-sense”\fnote{48}{(*)}{在“我那个年代”,英文里那个沿用已久的词“général non-sense”(泛泛之谈,意为:有时烦人却往往必要的通论)并没有贬义,而更多带着一点玩笑与憨直的意味。那个沿用已久的定语“général”在这里被“忘掉”,肯定不是偶然:这样一来就成了“non sensé”,用漂亮的法语说,不多不少正是“无意义”,还让人联想到“嗡嗡的空谈”(bombinage)和“胡扯”。}。他还有这份慈悲,不去细说SGA~4里肆虐的那份令人惋惜的non-sense(比如topos以及诸如此类的恐怖之物——读者可以庆幸自己福星高照:靠了这部光辉文集的天意降临,终于对那可惜的“矿渣”来了一场大扫除……)(注67$'$)(注67$_{1}$)。

我此刻把这一卷的导言与各章导言通读了一遍,把在我看来最明白地亮明底色的那些评语与自白抄录于此,另有两三句(路数是:枝蔓嘛,当然,不过“非常有趣”)在我看来主要是用来“把苦药裹上糖衣”的(而且糖衣果然顺利咽下)。比如,作者倒有这份诚实,开宗明义地说清楚:“若求完整的结果与详细的证明,SGA~4仍然不可或缺”。这一卷,无论其精神与动机多么暧昧,并不属于一场诈骗\fnote{49}{(**)}{(5月26日)不过,请看后天的“颠倒”一注(n\textsuperscript{o}~68$'$):在那里,我回到了这一印象——它被证明是仓促的。在反思的进程里,一场大面积的“SGA~4½--SGA~5”战役渐渐显露,它主要是为了Deligne的“收益”而做成,有我全部“上同调学家”学生的协助或默许。我所相信可以确认的那份“诚实”(依据导言第7行刚引用的那句话),在这里扮演的,是“对照线”的角色,用以障人眼目——是地道的“暂停!”风格。我的朋友从1968年起就在用这种风格(见“罐装的权与十二年的秘密”与“驱逐”两注,n\textsuperscript{o}~49与63)。另可参见“暂停!”与“中国皇帝的新衣”两注,n\textsuperscript{o}~77与77$'$。}式的战役。它的角色,在我看来,更近于一次试探——显然是结论性的:原来根本用不着这么多顾忌!

\dmark{}\origpage{252}有一种\textbf{荒谬性上的步步升级}(显然无人察觉!):从这一卷,到他筹备的那一卷(SGA~4½,以及LN~900)。在两处,人们都看到一个手段出众的人——一个生来就该去发现、去巡游、去探测浩瀚世界的人——执意去“重做”一位先行者的工作:先是我,然后是我的一位旧学生(Saavedra);而与此同时,他对这些先行者的工作——它们都是细心做出、并且深探到底的——拿不出任何本质的新东西。(他所能加进去的东西,总共恐怕二十页、三十页就写得下,依我看。)在第一个个案里,给出的理由还说得过去:让非专家的使用者得以“无泪地”亲近平展上同调\fnote{50}{(*)}{(6月10日)写这条注时,我才刚刚“登陆”,还没有感觉到“SGA~4½之役”的真正意味(及其与SGA~5坎坷命运的联系——对后者,我恰在那时才有了一次突然的预感)。后来我懂得:以那个骗人的名义SGA~4½发表的这堆五花八门的文本(见“颠倒”一注,n\textsuperscript{o}~68$'$),根本不是要充当SGA~4与SGA~5研讨班(它们构成我已刊数学著作的心脏)的一部“无泪”的普及读物,而是一次用它取代我已刊著作的行动:把我的著作呈现为一个多少有些含混的先声,把它自己呈现为关于平展上同调的真正杰作——仿佛这杰作当归功于Deligne。对这样一种冒名顶替的惊人表述(出自一支匿名的手笔),可在“SGA~4½”这次“试探”六年之后,“葬礼悼词(1)——或恭维”一注(n\textsuperscript{o}~104)里看到。},不必倚仗卷帙浩繁的SGA~4与SGA~5。(不过,人们头一回看到这位作者对常人世界如此体贴,竟压过了做数学的乐趣……)第二个个案里,这项工作实际上等于\textbf{照抄}Saavedra跟我做的学位论文!那篇论文本是一个完美的参考;其中一个陈述的证明有误、另一个陈述带有一条多余假设,这肯定不是重写整篇文章的理由。当然,对这样一件怪事,没有给出任何“理由”。

然而,我无需亲手捧起SGA~4½,就感到了这件表面上荒谬的事情的意味:Deligne在十年之后“重做”Saavedra的论文!它无疑同那件只稍逊荒谬的事的意味相同,而后者为它做了铺垫:Deligne(在十二年之后)对Grothendieck已刊著作的某一部分做一个“摘要”(边缘上还带几分居高临下)。恰恰是这一部分,倘若他还继续对代数簇的上同调(他始终无法摆脱它)感兴趣,他在任何情形下都不能装作可以没有。而Saavedra的论文,在所有已刊的、带着我的影响之印记的工作里,恰恰是他无论如何都绕不开的那一部——倘若他想把我发展的\dmark{}\origpage{253}MotifGalois群观念“记在自己名下”,并终于(十五年后!)开发利用这个显然至关重要的观念。经由先执笔SGA~4½,五年后又经由LN~900里那篇江河浩荡的Milne--Deligne文章(别名Saavedra),我的朋友沾沾自喜地给了自己一种虚幻的解脱感——解脱于他分明感到是痛苦负担的东西:不得不时时引证那个他正要取代、否认的人,哪怕仅仅引证某个引证他的人。

要抵达这一亲身的信念——关于这两个“荒谬”举动所共有的意味——我根本无需通读我这位多产朋友那(五十一项)出版物的全单;那份清单,十天前我才(头一回)收到。老实说,我甚至没有想到把我手里的四份抽印本再翻一遍\fnote{51}{(*)}{此外,还不算载于IHES《Publications Mathématiques》的那些工作——所长Nico~Kuiper,将近十五年来,一直好意地把它们寄给我。},为我相信知道的事找佐证。倘若今后我还会查阅我朋友的著作,那将是为了在里面找到别的东西——那些我在别处已经知道得足够多的东西之外的东西。那时我肯定会有学习美好数学的乐趣——那些数学,当年我曾有幸当面聆听、亲口从他那里学习,那乐趣还要更大!

\notehead{67$_{1}$}{}

\begingroup
\renewcommand{\labelenumi}{(\arabic{enumi})}
\begin{enumerate}
\item (6月14日)在SGA~4½里,我又记下了两桩微型的(细节上的)巧取。其一在“SGA~4,SGA~4½,SGA~5的引路线”(请欣赏这串富有暗示的排列!)里:作者写道(第2页),为了在平展上同调里确立“一个类似于凝聚对偶的对偶形式主义……Grothendieck利用了奇点消解与纯粹性猜想”,从而给人这样的印象:这一形式主义,最终是靠他Deligne本人才确立起来的——在(对许多应用已然足够)0维或1维正则概形上有限型概形的情形(见同一段)。他心里非常清楚:六种运算的形式主义(即整体对偶理论)是我不依赖任何“猜想”确立的;他所加的这条限制,只对双对偶定理(或“局部对偶”定理)才有依据——而那条定理,到头来,在SGA~5里(在Illusie的笔下)还改了姓,叫“Deligne定理”!

另一方面,第100页上有一节,题为“Nielsen--Wecken方法”——那正是我在代数几何里引入的方法,用来证明一条Nielsen--Wecken型的公式;这两位作者(在超越语境里)是用一种三角剖分技术证明它的,那技术在\dmark{}\origpage{254}代数语境里毫无用处。Deligne知道这一方法(连同Nielsen与Wecken二位先生的大名——那篇漂亮的德文文章,他根本无需去读!),是从我嘴里、在SGA~5那场“技术枝蔓”的研讨班上听来的——而SGA~4½正是要让人忘掉那个研讨班!在这一节里,既没有对SGA~5的暗示,也没有对我的暗示;关于这一方法的父权,读者可以自行选择:要么归于Nielsen--Wecken(倘若他消息极其不灵),要么归于这位光彩夺目又谦虚过人的作者。

有意思的是:整部书里,Verdier那“Woodshoie”式的证明(照录原书拼写)——关于一条迹公式,包括我所需要(为了Frobenius态射)的那个情形——只字未提。这一证明(显然已被人遗忘,让位给SGA~5里发展的那个更一般的方法)本是把我的L函数的上同调诠释完全立稳的缺失环节。显然,Deligne与Verdier之间有过(想必是心照不宣的)交易——Verdier把Weil猜想迹公式的信誉让给Deligne,作为交换,Verdier保有他在前一年(1976年)据为己有的SGA~5那一部分。(对此可参见“好的参考文献”一注,n\textsuperscript{o}~82.)另一项补偿:导出范畴与三角范畴的“第0状态”在SGA~4½里的发表——我的名字在那里同样缺席。况且四年之后,在Deligne笔下,代数几何里的平展对偶又得了“Verdier对偶”之名——Verdier这笔买卖做得不坏!(见“不义——或归来的意义”一注(n\textsuperscript{o}~75)末尾。)

\end{enumerate}
\endgroup

\subsection{格格不入者}

\notehead{67$'$}{}

(5月27日)\fnote{52}{(*)}{本注出自前一条注“白板”的一条脚注,是对它的补充,恰好一个月后的同一天写成。}所引的那些段落,连同围绕这部名叫SGA~4½的卓越文集之出版的全部情势,见证了我朋友的一种蓄意的嘲弄与轻蔑姿态——对着我的著作的核心部分:由SGA~4与SGA~5这两个血肉相连的研讨班构成的全体。从4月24日(见“同伙”一注,n\textsuperscript{o}~63$'''$)到5月18日(见“遗骸……”“……与身体”两注,n\textsuperscript{o}~88、89)的反思过程中显露出来的这些“情势”里,原初的SGA~5研讨班遭到劫掠——落实为1977年的那场屠杀式出版——绝非最轻的一项。(尤其见“屠杀”一注,n\textsuperscript{o}~87.)

我朋友这份蓄意的嘲弄,若想起以下事实,便取得全部意味:口头的SGA~5研讨班,正是年轻的Deligne\dmark{}\origpage{255}与概形、与上同调技术、尤其与对偶形式主义、与$\chi$-进上同调的初次接触——1965年,他21岁,在IHES登岸,目标十分明确:跟我学“代数几何”。正是在这场口头研讨班里,以及在两年前SGA~4的笔记里,他得以有幸第一时间学习那些直到今天仍主宰其著作的想法与技术\fnote{53}{(*)}{差不多同样的话,其实也适用于我的其他每一位上同调学家学生:Verdier、Illusie、Berthelot、Jouanolou——对此可参见“团结”一注及其后四注(n\textsuperscript{o}~85至89)。}。

“SGA~4½之役--SGA~5”、乃至我朋友皮埃尔对我个人之关系的这一本质的背景,显然在我写前一条注(“白板(1)”,n\textsuperscript{o}~67)时并不在场,在我之前关于葬礼的反思里也不在场。那个“年轻的Deligne”的记忆——他登上SGA~5研讨班,那时他还一切都等着学,而且确实(而且很快)学到了许多——只是在反思的最后阶段,仿佛违背我的本意,才浮上来。从我那个小小的数学“微观宇宙”里出现年轻的Deligne那一年起,我心中便立下一条蓄意的主意:不把他算进我的学生之列(仿佛把他算进去,就会对一位天资如此夺目的人有失谦逊);这一主意使我还一并低估了——或者不如说,直到最近这几星期才完全忽视——一个分明而具体的实在:人们通常用“师-生”这双重的称呼(而这称呼正是我一向拒斥的)所表达的那个实在\fnote{54}{(**)}{(6月14日)这一蓄意的主意,分明见于我四个月前在“耶稣与十二使徒”一注(n\textsuperscript{o}~19)里终于决定如何谈论他的那种方式——仿佛这样一来,我就在对那个乐于与我划清界限的人,违背了一项缄默或谦虚的义务……}。我曾乐于忘掉、忽视:确实有过某种从我到他身上的“传递”——某种对我、对他都具有巨大\textbf{价值}的东西,尽管对其价值所在,他与我的领会想必大不相同。在他与我之间那四年密切的数学接触中,我所传递的,是某种倾注了我自己最好部分的东西,一件由我的力量与我的爱滋养的东西——一件(我相信)我毫无保留地赠出、却不曾掂量、甚至或许不曾真正感到其分量的东西。

可以肯定,我所给予的,喂养了他身上那种求知的热情\dmark{}\origpage{256}——它与我身上跳动的那份热情同调——但也喂养了\textbf{别的东西};那东西我很久之后才感觉到,而且当时不曾把它与那桩已经发生、却乐于被我忽视的“传递”联系起来。换个说法:我所给予的,在另一个对我隐蔽的层面上,\textbf{也}被接受了——不是作为探测一个迷人的、永不枯竭的未知的工具,而是作为(起初)取而代之、(后来)确立支配、确立对他人毫不容情的“优势”的\textbf{工具}。

不去区分我朋友身上归于那个渴望发现的“孩子”的份额、与归于那个渴望取代、渴望支配(乃至碾压)的“老板”的份额,而只从表面看某些观念、技术、工具在一部著作中所占的份额——即便如此,最近六个星期里的一项意外发现仍然触目惊心:我朋友的著作,从我们相识那年起就振翅起飞,而它直到今天,仍在由我传递给他的东西滋养着。十五年前退出数学舞台时,我曾想象:我传给我这位“非学生”的朋友的那“一点点”(尽管我分明看到它在一次惊人的最初冲刺中的作用),将是第一个跳板,一次起飞将把他带得远远越过起点,使他\textbf{远离}我的著作、远离我这个人。而实际发生的却是:我的朋友直到今天,仍系在那起点上,\textbf{系}在那部他既要否认、要么交付嘲弄或遗忘、又要“利用”的著作上。这是对一个与父亲或母亲的冲突性纽带的典型个案:那根纽带,无限期地把那个乐于在自身培养这份冲突、而不肯冲向世界的人,拴在他注定要离开、要超越的人的轨道上……

今天我看清:经由那份蓄意把我的年轻朋友当作“格格不入者”、而不简单当作一个有幸比其他学生更有才具的学生来对待的主意,经由那份在他面前低估或忘记我所传递之物的分量(以及由此放进他年轻双手的\textbf{权力}……)的主意——经由我的这些态度,我在不知不觉中喂养了他身上的自负与冲突;而这两者,都对我隐蔽。与此同时,我入了一场双人游戏——或者不如说,是一场配合默契的双人游戏,我实在说不好是谁“开的头”(假如这问题有意义的话):我以“谦虚”自命,说我的年轻朋友实在太耀眼,不该做任何人的学生,说我所能给他的一点点实在不值一提——而他,则(早在我离去之前)与我这个人、与我的著作划清界限,在我的\dmark{}\origpage{257}纵容目光之下,否认那分明滋养过他的土壤。

只有在写本注的这一刻,我才终于看清这场游戏;对它的散漫知觉,不过是一两个星期前才在场的。我也看清:我身上的这份“谦虚”或“谦卑”,是假谦虚、假谦卑——是一种不单纯:不肯简简单单地把事情看成它们本来的样子。在这场游戏里,有对我的年轻朋友的纵容——播下的种子滋蔓了百倍!——还有更隐微的,对我自己的纵容:为一桩“特权的关系”、一桩非凡之至的关系垒起一个高台\fnote{55}{(*)}{可与5月10日的“高升”一注(n\textsuperscript{o}~63$'$)对读:在那里,我头一次知觉到,在我与我的朋友皮埃尔的那段关系里,有纵容这一配料。这份知觉始终零散而残缺,直到今天,才在本注“格格不入者”所进行的反思中明确起来。}。(正如一切不单纯,说到底,多半是一种对自己的纵容……)

\subsection{绿灯}

\notehead{68}{}

(4月27日)老实说,我从不曾反思过SGA~5研讨班那离奇坎坷的命运背后的意味。它1965/66年的口头进行,并未引起什么特别的困难;而由一批批(常常掉链子的)志愿者所做的成文工作,却拖了\textbf{十一年}\fnote{56}{(**)}{以我为口头报告所作的详细笔记为底本,把整个研讨班写定,对我而言本不过是几个月的工作。}!直到1976年,Illusie才终于把事情接过来,着手写定尚付阙如的部分,并出版全帙。今天——在这个研讨班之后将近二十年——我才头一回明白“这里面有东西要理解”。也许只有我一个人……

我冒出的第一个念头是:在研讨班那些多少积极的听众中间——他们也多少熟悉此前的SGA~1至SGA~4——必定发生过一种\textbf{饱和}现象:对着那滚滚而来、如一场无可抗辩的怒潮般的“Grothendieck学问”\fnote{57}{(***)}{这与我那班“多少有点被震懵了”的学生的印象相合;那一印象,见“教学的失败(2)——或创造与自负”一注(n\textsuperscript{o}~44$'$)所引的信。}。显然,某些执笔者信仰不足,他们想必不太明白这一切要走向哪里,我为什么偏偏要那么固执,整整一年,要把平展上同调的那些\dmark{}\origpage{258}本质形式性质翻来覆去、直到完全驾驭,连同与之相系的那一整套新观念的武库。尤其能说明那种普遍冷淡的是:研讨班的收官报告——提出未决问题与猜想(据我所知从未发表)的那一场——和梳理各语境中Euler--Poincaré型公式与Lefschetz公式的开篇报告,竟都没有留下痕迹。我不记得当时曾知觉到这种冷淡(甚至后来也没有,直到今天\fnote{58}{(*)}{(5月26日)正是在我多少重新“泡进”SGA~5研讨班的语境之后,我才想起一种不安的印象:那是我翻阅刚收到的、已出版的研讨班文本时产生的(那应当是1977年,出版之年)。那份“残缺”的印象(当时只停留在弥散、未成形的形态)主要地——也许甚至完全地(我当时没有花多少工夫细看,尽管那很值得!)——是由于开篇与收官两份报告的缺席,而尤其(我相信)是由于宣告这一缺席时的那份漫不经心,仿佛这是理所当然——何必费那份力气把它们收进去呢!我在某个层面上必定“感到了什么”,但直到这个月(将近七年之后!)我才费心让它浮上来、加以检视——在“屠杀”一注,以及紧随其后的“遗骸……”“……与身体”两注里。}),因为我一头扎在当时的任务里。

SGA~5的命运——它原初拥有与我的任何其他研讨班同样强的\textbf{统一性},却在随后的十一年不成文的岁月里渐渐被\textbf{拆解}(注68$'$)——本可以让我看到:我以那样的执拗所推进的那些大计划,我曾有几年的时间找到人手辅佐,却根本没有成为一项共同的事业,而始终只属于我个人。我的纲领在这里或那里唤起了些应景的合作,却从不曾在我当时任何一个学生身上变成一种观念之力——一种会促使他去做比他跟我做学位论文时更长程、视野更开阔的工作的力量;那篇学位论文在他生命中的主要作用,是让他学会他选定的数学家这门手艺。

依我看,唯一一个从整体上把握了(且不说据为己有)某种总体视野的人——超越就某类问题、或为发展某类工具而进行的特定“合作”的框架——是Deligne。这想必就是为什么,我(无须把事情明言)在他身上看到的,与其说是一个“学生”,不如说是一位\textbf{天造地设的继承人}。“继承人”这个词,比先浮现在我眼前的“接续者”更好地表达我想说的——后者会暗示一部受制于所受遗产的作品。相反,我把这份“遗产”感受为一个单纯的\textbf{馈赠}:我有能力做出它\dmark{}\origpage{259},以助一种个人视野的铺展;那视野还将从许多别的馈赠中汲取养分(事实上,早在我离去之前就已经如此),并且注定将毫不费力地超越先行与滋养它的一切。

回到SGA~5的悲惨命运:昨天掠过我心头的念头是,这一命运与Deligne对我个人、对我的著作之关系的暧昧,恐怕不无联系——尤其考虑到他那强大的数学人格,必定对我的全体学生施加过一种支配力\fnote{59}{(*)}{(4月28日)这一支配力的一个具体而雄辩的征兆是:SGA~5的出版,直到Deligne认为合适、向Illusie示意让他积极张罗的那一刻,才终于完成——也就是说,恰好在他自己需要把它当作其SGA~4½“摘要”的底本、以便取而代之的那个\textbf{时刻}。(对此可参见Illusie执笔的SGA~5导言的结尾。)这一下照亮了、并使前天我在“白板”一注(n\textsuperscript{o}~67)里还称之为“神秘的”那句声明取得全部意味:“SGA~4½的存在,不久将使人得以照原样出版SGA~5。”“照原样”在这里是一记尖刻,我恐怕是唯一(从前天起)感觉到并赏识其分量的人!(相对于原初研讨班,已刊版本代表的那场“拆解”,人们有目共睹。)}。他必定在这份研讨班笔记的坎坷命运中,在其被剥除那构成口头研讨班之统一与冲劲的一切之后,在内心深处得到了满足。然而反思之后,有一点很清楚:这些坎坷的第一性的、本质的原因,\textbf{不在}与会者中任何一个人的性情里。这一原因虽然我尚未清晰地辨明,但毫无疑问,它首先关乎我自己的\textbf{这个人},\textbf{也}关乎那些在65/66年间假装承担研讨班成文工作的人们。它肯定就在他们对我的关系之中,或者也许,在他们对我为他们所化身的某种做数学的方式(或某个纲领,或某种看待事情的视野)的关系之中。SGA~5的命运,如今在我眼前,是一种雄辩而执拗的\textbf{显影剂},显影某种我迄今不曾费心检视的东西——只因为我甚至不曾察觉它;直到此刻,我也只是隐约瞥见\fnote{60}{(**)}{(5月26日)这正是前面倒数第二条脚注里所说的那个“某种东西”;它在过去这几个星期的反思中、尤其从那个时刻(5月12日)起,终于浮出了水面——那天,自它1977年问世以来,我头一回费心看一看,“一场辉煌的研讨班”落到我的上同调学家学生们手里、经他们十一年后做成屠杀版之后,变成了什么。}。也许,这几行文字会促使这场集体厄运的某位当事人,把他对这件事的切身印象告诉我。

\pmark \origpage{260}然而,也许有一课——至少暂时地——我现在就能从SGA~5这一情节里吸取;这一情节先预演、随后例证了我所投身的那个著名“纲领”在我离去之后几乎全线的那种壮观的\textbf{停摆}。与我在六十年代那些兴高采烈的年月里或多或少相信的相反(那时我正为自己终于找到好心人辅佐而喜不自胜!),今天在我看来:经由顽强而缜密的工作使一个宏阔的个人视野取得肉身,本性上不可能是一场\textbf{集体}的冒险或事业。或者不如说:倘若真有“集体事业”,那也不是那种围绕同一个人进行十年、二十年(甚至三十年)的工作。只要这一视野注定要成为人人共享的遗产,它就会仅仅在需要的压力下、在此处或彼处,经由这个或那个人日复一日的工作而取得肉身——那个人对这位先行者,也许只知道名字(充其量如此!);而那位先行者的视野太过宏阔,凭他一己之臂,无法使之成形\fnote{61}{(*)}{(4月28日)也许,“我的一己之臂”本足以实现我在六十年代末所设想的那项宏阔的工作纲领——条件是,在随后的二三十年里,我让自己成为这一纲领的独家仆人。我今天感到庆幸的是,我没有走那条路——它本可能是我的路,而我如今清楚地看到它的陷阱与危险。}

\subsection{颠倒}

\notehead{68$'$}{}

(4月28日)作为这场拆解的一个例子(此外还有许多\fnote{62}{(**)}{(5月28日)我只是到5月12日的反思里,才下决心盘点这场“拆解”——见那则名字更贴切的注,“屠杀”(n\textsuperscript{o}~87)。}),我重新想起了SGA~5一场关键报告的命运:它最终由Deligne本人(我相信他从1965年起就揽下了这件事,为的是十一年后才“履行”自己的承诺……)依据我的口头报告写成,然后不经任何手续,直接收进了SGA~4½!说的是:与正则概形上的代数闭链相联系的上同调类的形式主义——它经由过渡到所论闭链之支撑上的“具支撑”上同调,而顺畅地展开。几乎与平展上同调中的一切构造一样(它们在许多别的语境里同样有用,并且在那里已成惯例),这一构造是我在五十年代末在凝聚上同调的框架里发展的(这里是Hodge上同调与De~Rham上同调;在“抽象”代数几何的框架里,它们头一回得到研究,是在我最早的几场Bourbaki报告之一里)。它如此\dmark{}\origpage{261}自然,以至于显然蕴含与上积的惯常相容性\fnote{63}{(*)}{(5月28日)凝聚框架里,见我的Bourbaki报告n\textsuperscript{o}~49(1957年5月),§40(原书如此)。在5月8日的“好的参考文献”一注(n\textsuperscript{o}~82)里,我发现:这些想法,连同我在同一场SGA~5研讨班里为闭链所配的同调类而发展的那些想法(以及许多别的),已被 J.L.~Verdier据为己有,对SGA~5研讨班的存在、对我的这个人,都不置一词。这场战役发生在1976年,即“SGA~4½之役”前一年(在我看来,它与后者紧密相连),而且是在1965/66年母研讨班SGA~5的全体往届听众与与会者的眼皮底下、在他们全都知道的情况下进行的。}。

写这些文字时,我发觉这场戏法(把这场关键报告偷运进SGA~4½)竟成就了这样一个光辉结果:Deligne——他明明参加了65/66年的SGA~5研讨班\fnote{64}{(**)}{(5月28日)而且,正是在那里,他头一回听说了他在那部海盗文集SGA~4½里陈述得如此光辉的那些事情!对此可参见昨天的“格格不入者”一注(n\textsuperscript{o}~67$'$)。相对于他朋友Verdier前一年的手法、以及他自己在别的场合使用过的手法,我的朋友在这里倒是守在公然劫掠的界限之内:因为他把我列为那场闭链报告的作者(诚然,附带那项光辉的成果:把我列为他的“合作者”),而且尚未装作全然不知道:平展上同调理论、迹公式等等里面,有我的一份。至于在这一方向上的决定性进展,可参见“葬礼悼词(1)——或恭维”一注(n\textsuperscript{o}~104)。}——\textbf{不在}封面我的“合作者”之列(这件事昨天翻阅已出版的Lecture~Notes n\textsuperscript{o}~589时就已打动了我);而我,反倒有资格(在研讨班十一年之后)以“Deligne的合作者”的身份露面。不得不说,这是一次相当绝妙的\textbf{颠倒}!在SGA~4½出版之时——我就这样在不知情中与之“合作”——我停止一切公开的数学活动已经七年;以至于我甚至从不曾过问那个可怜的SGA~5的出版,它对我来说,属于一段已被我抛在身后的过去……

(4月30日)至于SGA~5,它如今显出这样一副面孔:一堆有点五花八门、没头没尾的文本(头尾都在半路丢失了!),只有参照SGA~4½的文本才“站得住”。有一件事引人注目,而我恰在此刻才注意到:SGA~4½这个名字本身,就分明暗示这份文本\textbf{先于SGA~5},而后者只有参照前者才存在\fnote{65}{(***)}{(5月28日)关于这场把SGA~4½“强行塞进”SGA~4与SGA~5这两个不可分割的部分——它们构成我已刊著作的心脏——之间的“暴行”的更深意味,见“遗骸……”一注(n\textsuperscript{o}~88)。}。倘若这份文本的作者\dmark{}\origpage{262}的心思不那么暧昧\fnote{66}{(*)}{(5月28日)“暧昧的心思”这个说法,在这里分明是一个委婉语!},而他出于情感的缘由,执意要把他的“摘要”(“外加若干新结果”)塞进他曾扮演过角色的SGA丛书,那么顺理成章的名字,当然是SGA~5~1/5。

我在这里看到第二场戏法;它使我掂量出:Deligne在SGA~5的命运中所占的份额,比我三天前所想的还要重。这也使我收回前一天表达的感想——SGA~4½并不属于一场诈骗式的战役。显然没有人(从Illusie开始——他的诚意当然无可怀疑\fnote{67}{(**)}{顺便也该趁此机会,感谢Luc~Illusie:他以那样的细心与忘我,把某些陷于困境的报告的成文工作、以及“全包”的出版,进行到底——而且是在肯定谈不上鼓舞人心的条件下;其中,我的全然缺席,肯定不是最轻的一项!(5月26日)在n\textsuperscript{o}~84至89各注、尤其“屠杀”一注所继续的后续反思之光里,我对Illusie的这些感谢,取得了巨大的、未曾预料的喜剧维度——写这些文字时,我远没有预感到!诚然,我是顶着我身上的一种抗拒写下它们的;那抗拒尤其表现为(已在计划中的)致谢在注的“正文”里被“遗忘”,以致我不得不靠一条脚注来“补上”。这份抗拒,无疑出于那种不安:从我头一回亲手捧起这部名叫SGA~5的文集起,我就感到了它(而且我相信,直到最近几个星期,我再也没有捧起过它);那份不安,我在前一条注“信号”(原书此处题名如此)的一条脚注(今日5月26日所写)里谈过。这份疏忽,恰好例证了在禅观中、对当下正发生于自身之事保持警醒的留意的重要性。缺乏这份警醒,这里的反思就停留在禅观之下,停留在一个表浅的层面——而留意到那份抗拒,本会引导我去探它的根,并由此也更近地看一看那场美好的研讨班变成了什么(这件事,我两周之后才做)。})——察觉这场“战役”,这无疑要归功于我已领教过的那份“支配力”,我想,还有我朋友这个人身上的魅力:二者都使他高悬于一切怀疑之上!

\subsection{化圆为方}

\notehead{69}{}

(4月27日)十一二岁的时候,我正被拘留在里欧克罗集中营(芒德附近);我发现了用圆规作图的游戏,尤其为六瓣的玫瑰花饰心醉神迷——把圆周用圆规的开度在圆周上截取六次,分成六等份,恰好落回起点。这一实验的发现使我确信:\dmark{}\origpage{263}圆周的长,恰好等于半径的\textbf{六}倍。后来(我想是在芒德的中学,我最终还是上了学),我在教科书里看到,这个关系据说要复杂得多:人们有 $\ell=2\pi R$,其中 $\pi=3.14\ldots$;我坚信书错了,书的作者们(以及自古代以来想必走在他们前面的所有人!)肯定从来不曾做过这个极简单的图——它明明白白显示,人们干脆就有 $\pi=3$。典型的是,我发觉自己的错误(把弧的长与连接其两端点的弦的长混为一谈)是在我把我的惊讶——惊讶于先行者们的无知——向另一个人倾诉的时候(一位被拘留者,玛丽亚,她志愿给我上过几节数学和法语的私人课),恰在我准备向她演示为什么应当有 $\ell=6R$ 的一刻。

一个孩子对自身理解力所怀有的这份信心——宁可信任自己的能力,也不把在学校学到的、书里读到的东西照单全收——是一件珍贵的东西。然而它不断遭到周围人的挫抑。许多人会在我这里讲述的经验里,看到孩童僭妄的例子:它终究得在习得的知识面前低头——事实终于戳穿了某种可笑。但一如我亲身经历的这一情节,其中丝毫没有挫折感、没有可笑之感,而分明是又一次发现(在我草率地把 $\pi=3$ 这一错误公式当作解释的那次发现之后):一次错误的发现,以及同时,人们应当有 $\pi>3$ 的发现——因为显然,弧的长\textbf{大于}连接其两端点的弦的长。这条不等式,恰好又向着被驳回的 $\pi=3.14\ldots$ 那个公式的方向;于是它显得讲得通了,同时我那时想必隐约察觉:也许有一些并非那么愚蠢的人,曾经钻研过这个问题。此刻我的好奇心已经满足;我不记得当时还想对这个数——它想必重要到人们为它单配了一个字母——的来龙去脉知道得更多\fnote{68}{(*)}{(4月28日)前面的追述,唤起了别的回忆;它们表明,那个大名鼎鼎的 $\pi$,比我起初记得的更让我着迷。在一本书里(也许就是同一本)找到的近似值$344/133$(原书如此)打动了我——它那么漂亮,我简直难以相信它只是近似的!那时我除了分数之外不知道别的数,我为表达 $\pi$ 的那个既约分数的分子与分母会是什么模样而着迷——它们必定是些很不寻常的数!不必说,在这些关于化圆为方的孩童反思里,我没有走多远。}。

\dmark{}\origpage{264}这一经验,无疑是头几课之一:当我自己的理解似乎与公认的知识相抵触时,要有某种审慎——这样的局面可能值得仔细检视。审慎是经验的果实;它贴合并补足(而不改变)对自身认识与发现能力的自发信赖,以及认识到我们本就具有这种能力所带来的笃定。

\subsection{葬礼}

\notehead{70}{}

(4月28日)昨晚重新想到SGA~4½封面上那一幕——我在不知情中,以我显赫前学生的“合作者”的身份露面——这事在我看来难以置信到这种地步,以致我疑心自己的记忆背叛了我,怀疑我其实分明被征求过意见、并且没多想就点了头。但这一假定,与我直到去年还抱持的态度背道而驰到如此程度——即我根本无意再发表数学(更不用说以某人的“合作者”的身份,而且是某位与我的关系当时在我看来已带着深深的暧昧的人)——以致它比它本该“解释”的那件事还要“难以置信”得多;而那件事,归根结底,对我没有任何神秘或不可解释之处!为了求个心安,我还是查了我朋友1976年至今的信件(为数不多,很快查完),当然没有找到任何关于SGA~4½出版的暗示。我还是给当事人本人写了几行,请他对这出我不太欣赏的“闹剧”给我些解释……\fnote{69}{(*)}{(5月26日)我的朋友到底赏了我一份回答,它驱散了最后一丝疑云。他让我以“合作者”的身份露面,确确实实是因为那场SGA~5报告——他把它写成并收进了SGA~4½;而他不认为有必要为这次移交、或为以“合作者”的身份露面征得我的同意,也不认为有必要给我寄一册这本我曾“鼎力合作”的文集,鉴于“我已经七年不做数学了”。(6月5日)我刚收到(迟做总比不做好!)Contou-Carrère的一封信(日期为5月30日),回复我4月14日的一封信——我在信中问他(为求心安):在我的书里,他是否曾见过一份SGA~4½的副本。看来确有这样一份副本,Contou-Carrère留在了自己手里(除非是他买过而忘了?)。另一方面,Deligne的回答似乎确实证实:他不认为有必要寄我一份:“给你寄一份4½的副本,本确实可能是个好主意;我无疑料想你那时不会看出它的趣味”(5月15日的信)。}

三天前,当我在反思里提到三年前我与我朋友皮埃尔的关系中发生的那个转折——我失去了\textbf{兴趣}\dmark{}\origpage{265},不愿再同他交流数学问题(见“两个转折”一注,即注66)——我想起了一份当时强烈在场的印象。要把它安放好,我得先说明:在流逝的那十年里,当我的朋友事实上扮演着我唯一的数学对话者的角色时,我曾指望(一如我指望他扮演那个角色一样想当然):他会成为我向他吐露的数学反思与\textbf{想法}的中继站,把它们转达给可能感兴趣的数学家们。如我在别处解释过的(见第50节“一段过去的重量”),正是拥有这样一位对话者-中继站的感觉,赋予我那些断断续续的数学活动时期以比解馋更深的意味:把它们联系到一项超越我个人的共同冒险。无疑也是这份感觉,使得在那么长的时间里,我从未感到丝毫发表我所发现的东西的欲望,更从未感到丝毫对自己退出数学舞台的懊悔。(这样的懊悔,说到底,从未出现过;我是毫无蓄意地、“重新出现”在那个“舞台”上的,甚至在我自己察觉之前!)

至于我的朋友在多大程度上回应了这一期待,我说不好——很有可能,只要他对我还保持着那份数学上的热忱——它由好奇与亲切的同情共同推动,正是它使他在我与数学家世界的关系里(也在某种意义上,在我与数学本身的关系里)所扮演的那个特殊角色成为可能、并显得自然而然——他就一直扮演着那个角色。一两天前,当我向自己提出前面那个问题时,我收到了(仿佛作为部分的即答!)Larry~Breen的一封信,寄来1974与1975年间几宗通信的抄件,其中有Deligne1974年写给他的一张两行短笺,附着我刚写给Deligne的一封关于Picard叠形式主义的信的抄件,征询Breen对我那封信的意见。他在其中以“大师”一词指涉我;我从中听出一种半开玩笑、半亲切的声调。我不记得还有别的场合,有人把我1970年离去之后向我朋友吐露的事情,经他人之口回传给我。很可能有过,而我忘了;更不必说,即便在我数学活动的那些片段里,需要请教朋友的时候也相对少见;而且直到1977或1978年,我向他吐露的那些反思,分量都有限。所以,\dmark{}\origpage{266}严格说来,在那一刻之前,并没有多少东西可供“中继”\fnote{70}{(*)}{我可以算作例外的,是我关于分层结构的剖视理论的最初反思——七十年代初,我想必向Deligne提起过。他以一种纵容的同情接待了我在这方面的期望,有点像人们对一个什么都不怀疑的大孩子的态度。(他在与我的关系里常常是这种心绪,而且肯定常常有理!)我朋友的怀疑——它由他所熟知的某些野性现象的知识所推动,而那些现象是我所不知道的——当时并没有说服我;毋宁说,他向我指出的那些事实,使我从那时起就疑心:人们为“做拓扑”而通行的“拓扑空间”语境,不足以柔顺地表达某些我觉得本质的拓扑直觉,比如“管状邻域”。在其后十年里,我几乎没有机会回到这些反思,我的那些“疑点”被我淡忘了些许;它们经由我1981年12月--1982年1月的反思重新变得现实(并成为亲身的信念)——那场反思,是被为“Teichmüller塔”建立一套“剖视”理论的需要所激发的。(对此可参照《纲领概要》第5、6节。)(6月5日)作为另一个例外,我可以算上我关于虚拟相对概形与虚拟Motif(在一般基底概形之上)的反思——我记得曾向Deligne提起过。由于这些东西与他已决意埋葬的一部瑜伽(直到1982年的开棺出土)紧密相连,他没有表现出对我向他解释的那些想法有任何共鸣,这并不奇怪——那些想法当然使我心醉神迷;关于它们的几点指示,见注n\textsuperscript{o}~46$_{9}$。}。

事情在1977年起变化:那一年,自六十年代以来头一回,我强烈地“咬住”了一种异常丰饶的实体。那是我关于地图的反思的开端,以及(大约同时)顺藤摸瓜地,关于正多面体的一条新进路的开端(见《纲领概要》第3、4节)。也是从那一刻起,对我很清楚:我刚刚把手指按上去的那些事实,开启了无法预料的视野,其广度与深度,可与Motif观念诞生时所隐约瞥见的(以及后来远不止隐约瞥见的)相比。

奇怪的是,那一回,我竟然还是带着他将成为回声的期待去找我的朋友——尽管“Motif”这个名字周围已沉默了七八年,那沉默本已足够雄辩,足以教会我:我的期待是幻象!这份惊人的不察,恰好例证了我心中那条蓄意的主意(甚至在一两年前发现禅观之后):对我与数学、与数学家们的关系,一概不予留意——它们本该属于一段遥远而彻底翻篇的过去!我头一条向着这个\dmark{}\origpage{267}方向的反思\fnote{71}{(*)}{关于这场反思,见“扫兴的老板——或高压锅”(第43节)。}恰恰发生在1981年——那是我与我朋友关系中的第二个“转折”之年,我已有机会谈到它。但即便在那场延续了几个月的禅观里,与其他数学家的关系也只被轻轻拂过;而与其中那个无疑最亲近的人的关系(至少在我们共同热情的层面上)——据我所记,连拂都没有拂过。然而那本会非常有用!

但带着事后的认识、经由此刻的反思,有一点变得清楚:那一刻所发生的事——它曾那样使我惊讶、使我失望(一次隐约的轻蔑突然现身,而我本以为可以分享一次给了我深刻印象的发现那犹带体温的喜悦)——正是注定要发生的事。恰恰是我要传达的东西的\textbf{分量}——它曾激起我对一种与我的兴趣合调的兴趣的期待——应当在我朋友身上,头一回激发出\textbf{使我气馁}的反射。这一反射理应更加强烈,因为从SGA~4½出版那一刻起,我就已经被“预葬”了。三年后我再次进攻时——那时我的朋友(手持他关于绝对Hodge闭链的漂亮定理)正准备把下葬办得像模像样,随即有次年问世的“值得纪念的文集”\fnote{72}{(**)}{即Lecture~Notes~900那一卷;见“一个梦的回忆——或Motif的诞生”一注(n\textsuperscript{o}~51)。}——同一个反射又起了作用,只是粗暴得多。(这一情节当时断绝了我们数学层面的交流,但并没有“使我气馁”……)

在两种情形里,那份漠然显然是真诚的,一如它在别的情形里、对着我之外的人表达时那样。这已不是我头一回在他身上(或在别人身上)看到:与求知的渴望异质的力量,把那渴望中和掉,并取代数学家的嗅觉。

正是在这两次机缘——1978年,然后1981年——我头一回如电光石火般瞥见了我朋友身上那桩矛盾的“\textbf{代价}”;那桩矛盾,多年以来我就知道,但作为他著作中、他对数学事物理解中的一种阻碍与限制的分量,此前从不曾向我清楚显形。但只是在这一我已持续一个月的禅观——关于自我离去以来阴险地进行着的某场\textbf{下葬}——之中,这一\dmark{}\origpage{268}分量才终于渐渐显于全光之下。

在显见的层面上,这些天与这些星期里我所发现的下葬——几年前我已有预感,却从未想到把特定角色派给任何人——首先是对\textbf{我的数学著作}的下葬,经由它、并且首先,是对\textbf{我的这个人}的下葬。当然,在所有人中最有资格对这场下葬伸出援手(许多人在内心深处都在呼唤它)、并主持这场无名葬礼的,是那位不久前还在众人眼中俨然合法继承人的朋友。倘若他主持了葬礼,他肯定不是独自参加了葬礼!但更深一层:我的朋友在这些漫长的十二年间所悄悄下葬的,不是别人,正是\textbf{他自己};或者说,是他身上那样一种东西——那件不打动任何人、精微而难以捉摸、如一朵花或一枚果实的香气、无价的东西。($\Rightarrow$注71)

\subsection{墓}

\notehead{71}{}

但顺着联想的线索,我已离开本题——本题是要唤起某一记“强印象”;三天来,它的回忆执拗地回到我这里。这一印象,定位于我与朋友关系中那个“转折”的时刻:当我面对一种蓄意鄙夷姿态的征兆(既含蓄又粗暴地显明)——那些征兆使我在数学层面上结束了我们的关系。我当时懂得:时刻已到,延续这样一种关系,我再也无所期待;而“决定”自己做出了,没有分裂,没有懊悔,作为这份迟到的(而且非常不完全的)理解的第一个果实。

我心中没有愤怒,更没有苦涩。(我记得在我们交往期间,从未对我的朋友动过怒,也从未有过苦涩——除了我离开IHES的那桩情节的时刻;而且当时被包括在内的,也不止他一个。)但有一种悲伤:在同一个继续为我所珍的人的关系里,翻过这一页——尽管曾经把我和他系在一起的最强的纽带,已经干枯、已经死去。而作为此后几年里仍在的一根刺,还留着那份未被消化的挫折:我曾带来想要分享的喜悦,带给那个在我看来最亲近、最有资格分享它的人,却撞在一份自满的紧闭之门上。这份挫折,我看,最终是被我此刻继续着的禅观消解的。就在今天,它还在向我显示:发生在我身上的,正是注定要发生的;\dmark{}\origpage{269}而这份挫折的第一责任人,不是别人,正是我自己——是我自作主张,耽溺于某一实在的幻象图景,而不肯使用我健全的能力、用醒着的眼睛去看这实在。

正是在这份悲伤的底色上,也在这份期待的挫折的底色上,出现了那奇怪的印象——它当时不是作为一场反思(当时并没有发生)的果实或终点而来,而是作为一份直接的、无可置疑的直觉。那就是:我能在数学层面上对我朋友说的一切,以及多年来我已对他说过的一切,我是托付给了一座\textbf{墓}——或者已经托付给了一座墓。虽然我从不曾对任何人谈起这份印象,也从未在后来的反思里把它白纸黑字记下,我却清楚地记得:当时在场的,正是那座\textbf{墓}的意象,以及表达它的那个词(用法语),也就是我刚刚写下的那个词。这份“印象”或意象,想必在那一刻,作为某种理解的(不妨这样说的)视觉表达而涌现——那种理解,在某个层面上,想必早已成形、早已在场,是长年累月间必定发生过的一整组知觉的果实,尽管注意力不曾留住它们,记忆不曾登记它们;一些想必十分简单、十分分明的知觉,但我没有“留住”它们,因为它们对于我身上那位常常有权随意拣择的人……显得不合意。无论在当时还是后来,这份断然的意象,都未曾与任何具体的、可触摸的“事件”的记忆相连——一件朝着这意象方向去的事件,本可能在我心中唤起它;对这一骤然浮现的意象的记忆,此后想必也很少掠过我;今天是我头一回稍稍停留在它上面。

倘若当时没有浮现任何记忆或联想,那肯定是因为我连接待它们的最低限度的余裕都没有。奇怪的是,那时(如果我把时刻定位得准\fnote{73}{(*)}{(6月11日)一些旁证向我确认:确实如此。这个“第二个转折”,位于1981年的下半年。})我正投入一场关于我与数学之关系的禅观,而这一情节——它毕竟相当有力地,经由一个现在,向我诉说某段过去——竟没有使我想起中断反思的“线头”,好把方才发生之事的来龙去脉纳入反思;那件事,在我的一生里并非无关紧要。

\dmark{}\origpage{270}此刻(在唤起这一意象、并说出它一出现就与一切记忆或联想脱节……之后)浮现的第一个(老实说,也是唯一一个)联想,是我的Motif之“梦”所遭的命运——在数学诸视野之中最为我所珍的那一个。倘若那段过去,也许仍对我保有一种隐秘的支配,那正是经由这个梦——而这份隐秘的支配(我相信我在写下这些文字的时刻隐约瞥见它)本身,有着超越言词的、梦的力量。倘若在过往投资的遗产里——在一场对数学的满怀热情的投资的遗产里——曾在流逝的十年间出现一种未说出口而深刻的挫折,那正是看到:一片死一般的沉默,围绕着我视为活的东西,而我曾把它们当作活生生、精力充沛、只待冲向天光的东西,托付给我的朋友!我一走,能够而且有使命看护这次绽放的,是他,唯有他;把唯有他(与我一起)才亲密感到的东西,交到人人手中。而我虽然从不曾以这样或那样的字眼向自己说清这件事——甚至从不曾(据我所记)停下哪怕一念的工夫,想一想我留下之物的命运——但随着岁月流逝,在我身上的某个地方,我必定已经懂得:这个始终为我所珍的梦,我是把它托付给了一座“墓”。

于是,随着这一唤起、以及它在我身上激起的第一个联想,我看见一串别的联想循着它的尾流涌来,向我揭示:我分明刚刚触到一个命门——那也许正是头号要紧之处,我的数学家过去的(久被忽视的)重量由此施加。

但此刻,我想,还不是循着这些联想走的时候——反思的这一“最终”阶段已经拖得很长了。关于我的朋友皮埃尔,关于Motif,我想在这场反思里已经说得够多——肯定,对许多人的口味来说,已经太多!我相信,就这些注而言,是时候以一场对这场双重下葬之反思眼下教给我的东西的\textbf{盘点},来收束它们了。

\section{VI.~事物的回归——或全体一致的赞同}

\subsection{一脚踏进旋转木马}

\notehead{72}{}

(4月29日)\dotfill\fnote{74}{(*)}{我认为最好在这里免掉读者整整一页关于禅观的一般考虑——那是一种顾左右而言他的方式,是不肯进入正题的抵抗的征兆。}

\pmark \origpage{271}我觉得,就那桩占据着我的事情而言,应当做的描述与澄清的工作,就其“局部的图景”而言,主体已经完成。(况且显而易见:眼下这些注定要发表的注,只是实际工作的一幅缩影——在这里把汇成这一幅或那一幅“局部图景”的所有要素一一摆明,根本无从谈起。)当然,经由这同一项工作,某种总体图景也必定已经成形——还模糊,它等待着被表述,以便取得形态与生命,并向我说出它要对我说的话。自我昨天的反思以来,我感到它已蓄势待绽,推着我借给它一个声音。

老实说,昨天的反思(我刚刚重读)所主要教给我的,\textbf{不关乎任何别人,只关乎我自己}。我带着某种如释重负,看到反思回到对我自己的反思这一坚实土地——因为一个星期以来,它屡屡给我这样的感觉:牵连别人的这个人,多于牵连我自己。昨天的反思终于向我揭示了一件分明显而易见的事:即我对一段过去——我的“数学家的过去”——的执着之力,以及那部大名鼎鼎的Motif之“梦”在其中所扮演的特殊角色。

话一旦终于说出口,它的显而易见便跃然眼前——最新近、也许最清晰的征兆,便是那阵情绪:两年后才发现某桩“事件”之时所引爆的情绪,即Motif在我的前“学生”兼朋友的引导下,那场“悄悄的”(而且迟到的)“重返”数学兽栏!这阵情绪立刻表现为一场自以为已告终结的反思的重新开工——重新开工立刻落实为五十页回顾反思的洪流!于是(而这一确认,在这场不合时宜的重启中已数次来到我面前)似乎我还没有像一两个月前、在阶段性终点与一阵(绝非幻象的)解脱感的高涨中所以为的那样,“下了旋转木马”\dmark{}\origpage{272}——带着这样的教诲:“我并不比别人高明”,“学生超过老师,我无须惊讶”\fnote{75}{(*)}{见“打住这木马!”一节,n\textsuperscript{o}~41。}。然而这教诲并未阻止我惊讶——只要“学生”在一个我全然未曾预料的方向上超过我,就足够了!不过,倘若这教诲不曾阻止“我惊讶”,它在流逝的反思中却不止一次对我弥足珍贵,使我避开那些惯常的陷阱(或至少,避开这些陷阱中的\textbf{某一些})。

回到这份“支配”之力、这份我对Motif之梦的执着之力:这份执着在本书的许多别处已经显现——既在《收获与播种》里(其中Motif被不止一次谈起,而且用词已然雄辩之至),也在《纲领概要》里(在那里,“客观地”说本无须谈到Motif),或在《主题概要》里(在那里,Motif在一群活力充沛的雏鸡中间,有点像没有孵过的蛋)。在后一部文本里——它上溯到十二年之前,显然是怀着疏离的心境写就的——关于Motif的最后一段,据我看,是唯一能让人忽然感到一股暖意流过的段落……

引人注目的是:自我离去以来的这十四年里,这份执着从不曾向我显现;直到昨天,我才终于隐约瞥见这一显然的事实,直到今天才终于把它表述出来。在将近三年前的那场禅观里(1981年7月至12月),我终于确认了第一个显然的事实:即我身上对数学的那份热情的持存——它在流逝的岁月里,曾以十分雄辩的方式表达过自己。但我对一段过去的执着,据我所记,当时未被察觉,直到今天仍是如此。

然而,我必定是随着“一段过去的重量”那场反思才开始隐约瞥见它的——那场反思,是在对数学家过去的禅观似乎已经收官之时、出于求个心安而到来的(只除了一点:我还不曾知觉到这段过去的\textbf{重量}!)。而且我写它的时候分明感到,自己还停留在事情的表面,并未真正穿透。我随后不得不补加的那些注(先是注46、47)在当时把我引向一个方向——有好一阵,它使我离开我自己:先是把注意力系于一部数学著作(以及在我看来它最“重要”的那些方面),继而是那部著作的坎坷、他人在其中的角色,而不是我自己。

\dmark{}\origpage{273}我刚刚重读了“一段过去的重量”那场反思(第50节)。临近其结尾,我确实开始隐约瞥见:那股“翻转之力”(向着一种不止是偶尔的数学投入)可能是“对一段过去的执着”——对数学家的过去——所使;但更是对“最近十年的过去”、对“1970年以后”的过去,而不是对那些已白纸黑字写下的事物的过去、做成了的事物的过去、1970年以前的过去。几行之后,我却回忆起——但只是“顺带地”——在“我那时眼前的宏阔纲领……只有一小部分得到实现”。写下这几行时,我所想的,主要是“宏阔纲领”中那些直接可实现的部分;而它们的动机之力(!)却远不及“Motif之梦”所代表的那种。(它的论证——但绝不是它的表述——在当时,是“地平线上”的大任务之一……)

显然,我对“Motif之梦”的执着(一如想必一切执着),首先(即便不是完全)具有自我本位的性质。那是不止\textbf{贡献}于一项共同事业的欲望,也是看到这一贡献\textbf{得到承认}的欲望。假定“Motif的宏阔图景”确实已按我自六十年代末所见的那般幅度被勾勒出来,但我在这一视野的绽放中的那一份却被隐瞒——我的不快,想必不会小于(甚至可能大于?)我得知那部“值得纪念的文集”时所感到的(在其中,我分明看到被我提炼并带到光天化日之下的某些观念与想法被重新拾起,却(至少我是这样感到的)被剥去了那曾令我深深心醉的呼吸与炽烈的生命)\fnote{76}{(*)}{(6月14日)这份“不快”,在我看来,首先出于那份无耻之感:一种对某条纽带的蓄意蔑视,装作无视它、把它看作无足轻重。当一个人发现的观念或结果被他人重新发现时——这是常有的事——情形全然不同。}。只要那要看“得到承认”的自我本位欲望——它对着我的数学过去的这些或那些或远或近的事物——尚未烧尽,我自称“下了旋转木马”就还为时过早。数学的“木马”已不再\textbf{包容}我,一如它当年包容我、一如它包容我的某些朋友。但我肯定还在上面留着一只脚;而且我疑心,只要我还掺和做数学,这只脚就会留在上面!

\subsection{事物的回归(或一脚踩进菜盘)}

\notehead{73}{}

(4月30日)我方才重新想起了SGA~5研讨班的命运,以及这一命运与SGA~4½的出版相系的方式。一个曾经混乱的局面——我只是最近这些天、并以\dmark{}\origpage{274}顺带的一瞥检视过它——如今在我眼前显得非常清楚。我刚就此给三天前的反思补加了一条脚注\fnote{77}{(*)}{这条长得犯规的脚注,变成了一则独立的注:“颠倒”(n\textsuperscript{o}~68$'$)。}(见“信号”一注,即注68(原书此处题名如此)),而且我觉得,加上我前天已有的评注(同样以脚注的形式)与此前一天的反思(“白板”,注67),我已表述得足够清楚,无须再为如今已显出足够雄辩的局面再做一幅总结的全景\fnote{78}{(**)}{然而5月9日及其后几天,我还要回到它;见注n\textsuperscript{o}~84--89。}。

行文至此,有一点很重要,必须确认:SGA~5所遭“悲惨命运”的第一责任人、那份“被弃局面”之利用的第一责任人,不是别人,正是我自己。倘若各位“志愿者”(他们揽下了自己并不真想做的成文工作)显然对自己心里没数,那我也同样没有看清自己:我执意不听从一个已然雄辩的局面给出的教训,把事情托付给三心二意的“合作者”,而不肯亲自接手、亲自去做那从那时起就落在我头上的成文工作。毕竟,从口头研讨班结束,到我离开数学世界的那一刻,整整三年过去了(而后者立刻表现为:在其后的十四年里,我对我已刊著作几乎全然的漠然)。诚然,那三年里我完全被别的任务占据:SGA研讨班的延续(SGA~6与SGA~7),EGA的写作,对日复一日涌现的、常常多汁的问题的反思——而在这些问题之中,Motif总体视野的渐渐成熟……被这些任务缠住,我选择了对一场已成过去的研讨班的命运闭上眼睛——那场研讨班(连同前一年的SGA~4),是我就已完成的工作而言所能做出的最深刻的数学贡献,而且多半也是影响最广的一份。

在我一去不返地离开之后,局面只可能更加恶化,使我那些前学生中最显赫的一位得逞:他把那部大名鼎鼎的SGA~4½,插进SGA~4与SGA~5那non-sense与多余细节的矿渣之间,上演了那出绝妙的战役;并让我荣任那份以中心关键文本自居之物的“合作者”——那文本的使命(如他以那种使他生辉的天真所说)是慈悲地让人“忘掉”包裹着它的沉重矿渣……

\dmark{}\origpage{275}总之,我早在离去之前、并经由离去本身所做的种种选择,蕴含着对我已刊著作(或就SGA~5而言,待刊著作)命运的影响,同样也影响我“著作”中仍停留在梦的形态的那一部分——而且是\textbf{未发表}的梦。我不后悔我的选择;当我今天确认这些选择的某些不合我意的后果时,也轮不到我抱怨!倒是轮到我去检视这些后果(而且越不合我意,越要检视!),对事实\fnote{79}{(*)}{(5月28日)这里应读作“我所知道的那些事实”。从后天起,一些全然意外的新事实,将使关于葬礼的反思重新启动,并使我把它相关的注的篇幅增至三倍。}做一幅总体图景(这件事已经做成),并从中汲取它们能给我的教益。这是留给我要做的;今天的反思,至少,也许是在这个方向上的第一步。最近这几天里,我身上已做成了某些比照,我想首先把它们白纸黑字写下来。

在六十年代初的那段时期,推动我一般地在我的学生们身上投资的主动力,是找到“\textbf{人手}”来实现我的本能指认为紧迫而重要的“\textbf{任务}”(至少在我所特有的数学眼光里)的欲望。这份“重要性”肯定不是纯主观的,不是一个“趣味与颜色”的问题;而且常常(我相信),把我的某项任务当作己任的学生,分明感到它“够分量”,也许还感到它在更宏阔的蓝图中的位置。

然而,就这份推动力、这份催促我实现任务的内在动力而言,所关的并不是某种“客观的”重要性——Fermat猜想、Riemann假设或Poincaré猜想的“重要性”,曾使我全然无感,我并不真正“感到”它们。把这些任务与一切其他任务区分开来的,在我与它们的关系里,是它们是\textbf{我的}任务;是我感到了它们,并把它们变成我自己的。我很清楚:感到它们,曾是一项精微而深刻的工作、一项创造性的工作的终点,它曾使人得以把握作为某项任务之对象的关键观念与问题。它们曾经是——而且无疑(在很大程度上)至今仍是我这个人的一部分。把我和它们系在一起(或今天仍系在一起)的那条纽带,在我把某项任务托付给一个学生时,绝没有被剪断——恰恰相反,这条纽带获得了新的生命、新的活力!这条纽带无须说出\dmark{}\origpage{276}(而我在这里把它“说出”,哪怕只对自己,也是头一回)。这条纽带,对那个选择了与我一道工作、并就其自选的某项任务而工作的学生来说,是明摆着的;对我,以及(我确信)对任何别人,也一样。这是构想了一件事的人与这件事之间深深的纽带——它不被后来那些也把这件事变成“他们的”、并向它贡献出自己最好部分的人所削弱,而(依我看)恰恰被加强。这条纽带,我从不曾细心检视过。它在我看来,深深扎根于“自我”的本性之中,具有普遍的性质。人们有时装作无视这条纽带,仿佛自己高踞于这类小家子气之上——甚至可能,我也曾一度入过这种装模作样\fnote{80}{(*)}{可以肯定的是:我遵循着那种“得体的做派”——即对这类事情视而不见,因为它们有悖于时兴的形象!(5月30日)关于这条纽带,见“……与身体”一注,n\textsuperscript{o}~89。}。但最近这些年里(或最近这些天与这些星期里),那寥寥几次:当他人摆出无视这条纽带——他明明知道这条纽带——的态度,而这条纽带把我系于某项已由他人(或由我自己)完成或仅仅指定的任务时,我被触到了痛处。人们可以把这个痛处叫作“虚荣”或“自负”,再给它扣上别的字眼——我并不声称这些字眼在这里不合用;但无论人们给它什么名字,我谈论它毫无羞耻,也不因自己如今的样子而羞耻;而且我知道,我所谈的这件事,是天底下最普遍的事!当然,一个人对其“作品”的这份执着,其力量并不千人一律。在我的生活里——“做”自童年以来就是我几番重大精力投资的恒常焦点——这条纽带曾是强的,今天依然是强的。

于是我可以这样说:推动我与学生之间关系的主动力,是我把他们看作可欢迎的“人手”,用以实现“我的”任务。这一表述可以显得犬儒,而它只是表达一个明摆着的实在——我的学生们想必同我一样分明感到了它。这些任务是“我的”,丝毫不妨碍他们也把它们变成“他们的”——正是他们在自己身上对任务的这种认同,动员了完成任务所需的能量;一如对这个任务的认同,曾在我身上动员了使它诞生并成形的能量,并且继续动员我持续投注于这一主题的能量。这份能量是不可或缺的,否则我甚至无法作为“老师”而“运转”——也就是作为把一门手艺(它同时也是一门艺术)传授下去的长者;而这件事,离开可观的能量的动员,无从进行。在我做教师的全部过去里,我从不曾感到矛盾\dmark{}\origpage{277}:同一个任务,对我这里工作的学生而言深刻地是“他的”,同时又深刻地保持为“我的”。我不认为这一局面带丝毫冲突的性质,也不认为它曾给冲突的意图提供过攀附的机会\fnote{81}{(*)}{倘若,受到某种气候的怂恿,我的一个学生曾想在某项与我合作的工作里藏掖我曾有过的角色,那么事情发生在一个他早已不在学生之位的时候。}。在这种对同一任务同时投资、并认同于它的局面里,学生和我自己(依我看)都各得其所,处于一种全然清晰的工作关系之中;而就其自身而言(依我看依然如此),它不含有任何冲突性的成分。但在纯属个人的层面上,这份关系却停留在表面——这丝毫不妨碍它亲切,甚至友好,有时甚至充满感情。

对我的任务的投资,并\textbf{经由它们}对我的学生们-合作者们的投资,是(我说过)自我本位性质的(想必一如一切投资)。这些任务的实现,对“自我”来说,当然首先是一种自我扩张的手段:经由实现一项宏阔规模的总工程——“我的一己之臂”本无法完成到底的工程。从我数学生涯的某个时刻起,就有这份恒常的暧昧:一种同居、一种紧密的相互渗透,一边是“\textbf{孩子}”——他的求知与发现的渴望,他对隐约瞥见之事、对细察之事的惊叹;另一边是\textbf{自我}、那个“\textbf{老板}”——为自己的作品而欣喜,渴望扩张,渴望经由作品的倍增、或经由一项宏大规模的总营造的执拗而不懈的追求,来增添自己的荣耀!在这份暧昧里,我看到一种分裂,它继续压在我的生活上,给它打上深深的印记——一种也许将伴我终生的分裂。这样的分裂,当然不为我这个人所独有;但也许,在我那被“最好”与“最坏”一并充满的生活里,这份分裂采取了比在别人身上更极端的形式。

于是我可以这样说:对那个贪得无厌、渴望扩张的“自我”而言(它在场上的并非孤身,但确实在!),我的学生们首先是可欢迎的“合作者”,甚至不妨说“工具”——可欢迎的“人手”,用以营造一座会诉说“我的”荣耀的巍峨工程!\fnote{82}{(**)}{写下这句话时,我略有踌躇,字斟句酌,明知有人会把它夺走,当作那个可怕的学阀终于撕下面具的犬儒自白!但我也知道:我拦不住那个一心想淹死一条碍事的鱼的人,他自会从容行事。这不会阻止我继续我的话题:发现并说出明摆着的事情——包括上面写下的那份寒微的真相;它只会让那种从不曾费心朝自己内部看一眼的人感到意外。}

在我看来,这一点已经相当清楚地\dmark{}\origpage{278}在三年前那场关于我与数学之关系的禅观里(进而,关于一般的“做”)显形过,尽管后来我也许把它淡忘了些许。正是这件事,在这最后几天里在我的思想中在场,使我得以同另一件值得注意的事实比照:恰恰是我那时的一位学生(带引号的学生,姑且如此!);而且在所有人中与我最亲近的、也是唯一能不费力气地、就其整体地“感到”我身上那些宏阔蓝图的人——那些蓝图仿佛不容分说地催促我去实现它们——恰恰是他,在我离去之后(而且在内心深处,想必更早……)经年累月地实施了那场\textbf{葬礼}(Enterrement)——规模与著作(Œuvre)本身相当的葬礼(这两个法文词在这里用大写并不过分!),并且最终“主持了葬仪”(Obsèques,再多一个大写,好压得住秤!)。

在这份局面里,打动人的,是那种乌布王式的、巨大的、不可抗拒的\textbf{喜剧性}!最近几天里,我想必隐约感到了这份喜剧;但它只是在刚刚,在我为我自己那场庄严的葬礼放上最后一个大写的瞬间,才向我显出它的真面目——在一阵突然而不可抗拒的大笑里!恰恰是这份\textbf{笑},在反思的这个号称“最终”的阶段里,直到此刻一直缺席;在这里,主导的音色,宁可说是那位“体面先生”的愁容——他的正当期望落了空(甚至可以说被卑劣地欺骗),当愁容没有让位于那些抖得干净利落的嘲讽评点的时候(会说话的人自然会说话!)。我分明感到,经过这场长时间的、忧伤音调的枝蔓(这个词让我想起了什么……),我重新走上了正路。

而此刻,一个顺理成章的名字浮现出来,用以命名这条“注”(人们已不再清楚一条注是关于什么的注,但无所谓),是收束它的时候了。它将是:“\textbf{事物的回归}”。($\Rightarrow$注74)

\subsection{全体一致的赞同}

\notehead{74}{}

我终于感到——呼!——我触到了这一“最终阶段”的终点;它绵延了十二天,其中每一天都(一如从前)以“最后一天”的面目出现。也许,终场的话,就在几分钟前已经说出。我的(象征性的)下葬,是一场\textbf{事物的回归},一场由我亲手播种的收获。(而我那有血有肉\dmark{}\origpage{279}的下葬——倘若我有这份幸福死去,身后留下能为我送葬的活着的男男女女——也将是一次回归:回到我出生时就离开的某物之中……\fnote{83}{(*)}{(5月28日)这一与我自己的死亡的突然联想,带着力量来到了。我一度想把它推开,随后又想删去这突如其来的括号——它仿佛饭里的头发,来得不是地方。我克制住了,出于某种敬意。奇怪的是,第二天我得知:就在我继续反思的那个4月30日晚上,在我居住的市镇,一位朋友的重病缠身的姊妹去世了。我头一回见到德尼丝,就在当天,在她的灵床上。次日5月2日,我与我的朋友、与许多别的活着的男男女女一起,在一个明媚的春日里,把她送入土中……})。我觉得,还可能要补加的,多半只会是\textbf{尾声}方面的了。

那位大名鼎鼎的“最亲爱的学生”,并不是我亲爱的学生们中间唯一一个兴冲冲为我送葬的人;而那些分明伸过手的人们,在他们中间,恐怕也不是唯一出席葬礼而不以为苦的!但说到底,谁是张三谁是李四,于我何干!(关于这件事知道得更多——如果只是为此——不会教给我任何新东西。)我终于懂得了这场“事物的回归”;懂了,我也就领受了它的益处。

然而,我还没有取尽这份益处为我储备的全部实质。我还不清楚地辨明:究竟是我身上的\textbf{哪样东西},使得某些前学生觉得,下葬与葬礼于他们有利可图。难道仅仅是我谈过的那份“贪求”?它(依我看)与别的“老板”们的贪求相去不远;而当年他们与我初试身手时,曾不费力地与它相安(而且无疑甚至不曾察觉,至少在意识的层面上)?那么,是“机缘”(我的离去等等)“造就了贼”,并且\textbf{显影}了他们身上——一如在“众学生之首”身上——的一种\textbf{普遍的倾向}:一旦情势有利,就埋葬自己的“老师”或“父亲”?也许,我的“老师”(或“父亲”……)形象还超出了通常的尺度,而这一情势,曾有助于以完美的齐奏引爆这套“下葬综合征”?此刻我看不清!也许,我将收集到的回声(但愿如此)会使我看得更清,更好地消化此刻摆在我面前桌上的这份意外的食粮。

参加下葬与葬礼的,并非只有学生们,尽管据我所知,没有一个非前学生在其中扮演过突出的角色。显然,我的许多旧朋友从中各得其所。这回,这件事在我看来并不太神秘。

\dmark{}\origpage{280}如我已有机会顺带说过的,我不止一次确认:我不合时宜地退出数学舞台,在我昔日的朋友们身上造成了深深的困扰。那是这样一类东西所激起的困扰:人们在其中隐约感到一种对深刻的重新审视、对革新的\textbf{挑衅}。在这个案子里,自然,数学家中间的这份困扰在我的朋友们身上最烈——在那些认识我、并能感到我曾对那些至今仍是他们的价值投入过多大力量的人中间;更不必说,这些朋友中的每一位,自己也曾在这些价值中、并在这些价值所提供的丰厚“回报”中,做出并且继续做出一份力量相当的投资。早在“生存与生活”(Survivre et Vivre)时期一开始,我就已有充分的机会在别的科学工作者中间观察到这样的困扰。但这从未妨碍:每一次,当我在某位昔日朋友的身上——同情的纽带依然把我和他系在一起——确认出疏远的、有时是敌意的毫不含糊的征兆时,我都感到意外。使我的“弃绝”对某些人变得格外难以容忍的,恰恰是:我本该是他们中间的“最好者”之一,是他们最不可能疑心会给他们来这一手的人!(而且,我确实好几次以为感到,数学界的某位旧朋友身上有一丝\textbf{怨恨}的音调。)于是,很自然,他们在一场时尚里得到了满足:那场时尚宣告,所有这些“Grothendieck学问”,说到底不过是一大摞纸换不来多少东西,等等等等。一个人,无论多么显赫,不足以成一场时尚——要发起的时尚,还得回应许多别人心里的期待、暗藏的欲望,然后才成为共识、才取得支配地位\fnote{84}{(*)}{(5月28日)同一方向上,可参见5月14日的“掘墓人——或全体教众”一注,n\textsuperscript{o}~97。}。

在我离去之后的这十四年里,我也许一直倾向于低估这一离去在“上流社会”里造成的困扰——而对我而言,1970年6月的这次离去进行得如此自然,以至连一项“决定”都无须做出:新的任务一夜之间接了旧任务的班;旧任务忽然退后,仿佛被一段遥远的过去吸收殆尽!(诚然,在蒙彼利埃大学的同事们中间,我也不曾面对这样的困扰;他们构成一个与我离开的那个环境全然不同的世界。)也许,我同样低估了这份困扰在我的“1970年以前”的前学生们中间所能起的作用——他们中的许多属于同一个环境,并在其数学投资中“倾其所有”。很有可能\dmark{}\origpage{281},这份困扰在他们身上起的作用,不亚于在我以为还拥有于同一环境中的别的朋友们身上。无论如何,(我的某位旧朋友或学生与我之间的)每一种情形,都是独一无二的、不同于所有其他情形的个案;我能做的一般推测,只有非常有限而临时的分量。

回到更坚实的个案的地面,有一件事打动了我:我得以确认曾积极参与亲爱的老师下葬的那两位前学生,恰恰也正是最初以轻蔑的姿态、以使人气馁的意愿引起我注意的那两位——对着更年轻的数学家们,他们是“1970年以后的学生”,或者我的观念与数学进路的影响在他们身上清晰可见。这一巧合当然毫不奇怪(这当然不妨碍事件每一次都使我惊讶!)。另一桩有趣的巧合是:两人都是个人关系最友好、甚至最充满感情的那一类(对其中一位,这份关系以同样的音调一直延续到今天)。这与那条一般的确认相合:恰恰是最亲近的关系,最有力吸引并固定冲突的力量。

还有一桩巧合打动了我。在我将近二十五年来有过的学生中间,有两位,在我看来,既以其出众的“手段”、又以其与这手段相称的数学投资,区别于所有其他学生。(一项与我在我生命的二十五年里所做的那种力量相当的投资。)而且,对两位,我都曾有所顾忌,不肯把他们算进我的学生之列——尽管确实,他们都在与我的接触中学到了对他们有用的东西\fnote{85}{(*)}{(5月28日)这是一个委婉语,正如我后来身不由己地确认的那样!对此可参见昨天的“格格不入者”一注,n\textsuperscript{o}~67$'$。}。事情的本性是:两人都各自发现了他们自己的任务,无须我从我曾有(或仍有)的储备中向他们提出任务;两人的学位论文工作,都是在与我这个人无涉的情况下完成的\fnote{86}{(**)}{(5月28日)这并不完全准确。两人都在其工作中,本质上使用了我锻造、并且是他们在与我的接触中学会使用的工具。在这一角色之外,构成其学位论文的工作里的Hodge--Deligne理论(“Hodge理论II”,《Publications Mathématiques》n\textsuperscript{o}~40, 1972, 第5--57页),直接出自他得自我的Motif瑜伽——“混合Hodge结构”是对(在Motif的眼光里同样“明摆着的”)问题——“把复数域上不必半单的Motif观念‘翻译’成(适当意义上的)Hodge‘结构’”的“明摆着的”回答。在一次出色完成的“翻译练习”之外,这项工作中当然有“与我这个人无关”的原创而深刻的想法。但同样清楚的是:倘若没有我引入数学、并且他们在与我的接触中率先获得的那些想法与工具,那套Hodge--Deligne理论今天就无从存在(而且恐怕Deligne或我的另一个学生的几乎全部著作,也无从存在)。}。如此多的共同点\dmark{}\origpage{282}!作为相异点,我要说:两人中较年轻的(如果我没弄错)今天“荣誉等身”(其详尽清单,我替读者、也替当事人那份众所周知的谦虚省去了),而且是影响力最大的数学家之一,也就是说,也是最有权势者之一;另一位,眼下是代理助理,暂代一个职位,而该职位的正式任职者明年就要回来。还有别的相异点,它们在某种程度上解释了这一命运的悬殊——一如也有别的相似点,在这里不值得一一铺陈。倘还要举一条,那就是:在我有过的全部学生中,个人关系最亲近、最友好的,恰恰也是与这两人;而共同的热情,当时一下子就在他们每一位与我之间建立了强的纽带。我此刻想说的这份\textbf{巧合}是:据我所知,他们也是仅有的两位学生(带引号,自不待言!),对着那个“上流社会”,曾尽其可能地淡化和抹去这条通向我的、如此简单而分明的纽带。

这是一份真正触目的巧合;写下这些文字的时刻,它的意味仍从我这里滑脱。对两人,我都可以援引各自的时势缘由,彼此不同。而且很有可能,甚至可以说大概:在两人身上,在某个大概已不再是充分自觉的意图的层面上,这样一种缘由(一方的自负,另一方的谨慎)起过作用。但我怀疑,这个现成的解释,无论在这个个案还是那个个案里,能提供对事情的理解。肯定,在更深处,别的力量——那些真正的力量——曾在自负或怯懦的熟悉表象背后起过作用。肯定,表达这些力量的行为,有某种重要的东西要向他们每一位诉说。但肯定同样:在两个如此不同的人身上出现同样的行为——仿佛他们串通好了(鉴于二人命运的悬殊,这当然不可思议)——也有某种重要的东西要向我诉说,而且不关乎别人,恰恰关乎我自己。难道又恰恰是那没完没了的“\textbf{对父亲的拒斥}”的再现?然而,这种拒斥\dmark{}\origpage{283}本有许多可以选择的表达途径!抑或,正因为无意识那如此可靠的本能——它(当涉及“触碰”时)总能不偏不倚地直取最敏感、最不设防之处——竟使得两人不约而同地落在了\textbf{同一个}地方?实际上,我倾向于这样想。但这是一件推论出来的事,而不是一件看见的事;而由于缺乏那双有幸看得清楚而深的眼睛,我感到自己有点像个盲人,在黑暗中尽力摸索,尽力用他的手、他的耳朵、他的皮肤去“看”——可它们生来不是用来看的。

然而,为了不以这条\textbf{困惑}的音调收束(它有损于我的名声),而以一条令一位善意的、或许会有的读者欢欣的音调收束,我将只说出方才浮现的那个一锤定音的名字;在我看来,它分明表达了这场(关于一场下葬的反思之)\textbf{尾声}的各条考虑所共有的内容,即:

\begin{center}\textbf{全体一致的赞同!}\end{center}
